% ============================================================
% 微模块4：人工智能赋能医学

% ============================================================

\part*{微模块4\quad 人工智能赋能医学}
\addcontentsline{toc}{part}{微模块4\texorpdfstring{\quad}{ }人工智能赋能医学}

% ============================================================
% 第9章：医学人工智能的应用基础知识
% ============================================================

\setcounter{chapter}{8}
\chapter{医学人工智能的应用基础知识}

在医院的诊疗过程中，一名患者从挂号、问诊、检查检验到治疗随访，会不断产生不同形式的医学数据。电子病历记录主诉、病史和诊疗意见，医学影像呈现器官结构和病灶形态，生理信号反映心率、血压和血氧等动态变化，生命组学数据则从分子层面揭示疾病相关特征。这些数据来源不同、格式不同，记录时间和医学含义也不完全一致，只有经过规范处理，才能用于疾病识别、风险评估和诊疗辅助。

医学人工智能常用的技术方法包括ML（machine learning，机器学习）、DL（deep learning，深度学习）和NLP（natural language processing，自然语言处理）。机器学习是使计算机从数据中学习规律的方法，深度学习以多层神经网络为核心，适合处理图像、语音和复杂模式识别任务，自然语言处理则用于识别和分析病历、报告及医嘱等文本信息。这些方法能够从医学数据中提取与疾病表现、健康状态和诊疗过程有关的特征，为病灶识别、疾病分类、风险预测和辅助决策提供支持。

医学人工智能应用需要同时关注数据质量、模型构建和临床使用过程。模型是基于训练数据形成的计算规则与参数体系，可在新的医学数据中识别与疾病状态、风险变化或诊疗需求相关的特征，并输出分类、预测或异常提示等分析结果。数据采集方式、标注质量、样本差异、图像质量和设备变化都会影响模型表现，因此模型输出不能脱离原始数据和患者临床背景进行解释。

医学人工智能的价值在于提高医学信息处理效率，辅助医生发现数据中的异常线索和潜在风险，但其结果并不直接等同于临床诊断结论。实际应用中，医生仍需结合患者病史、体征、检查结果和随访变化，对模型提示进行核对、修正和取舍。本章内容包括医学数据处理与任务类型、医学影像识别与辅助诊断、临床决策分析与健康管理，重点说明医学数据如何形成和处理，人工智能如何完成不同医学任务，以及模型结果进入实际诊疗流程时需要遵守的使用边界。


\clearpage
\section{医学数据处理与任务类型}

\begin{sectionkeypoints}
\item 了解电子病历、医学影像、生理信号和生命组学数据等常见医学数据类型，认识不同数据在信息来源、结构形式和医学用途上的差异。
\item 理解医学数据多源异构、专业性强、动态变化明显和敏感性高等基本特征，能够结合具体诊疗场景判断数据含义。
\item 掌握医学数据采集、清洗、标注和标准化的基本过程，理解这些处理环节对模型训练、结果分析和临床解释的影响。
\item 理解多源医学数据融合与表示学习的作用，认识患者身份、时间顺序、空间位置和医学语义对齐的重要性。
\item 熟悉病灶识别、辅助诊断、风险预测和治疗决策支持等医疗人工智能典型任务，理解不同任务与医学数据之间的对应关系。
\end{sectionkeypoints}


在医院的诊疗过程中，患者从挂号就诊、检查检验、住院治疗到出院随访的每个环节，都会产生不同形式的医学数据。患者的主诉、现病史和诊疗意见记录在电子病历中，血常规、生化指标等检验结果进入实验室信息系统，CT（computed tomography，计算机体层成像）和MRI（magnetic resonance imaging，磁共振成像）等影像资料存储于医学影像平台，监护设备则连续记录心率、血压和血氧等生理信号；患者出院后，随访记录、用药情况和康复数据还会继续反映其健康状态变化。

原始医学数据常包含记录缺失、单位不一致、重复检查、时间戳错误和标注差异等问题。如果这些问题直接进入模型训练，系统可能学习到错误关联，导致结果偏移，需要经过采集、清洗、标注、整合和质量控制等预处理步骤，才能用于后续分析。人工智能可根据数据特点完成分类、检测、分割、预测和生成等任务，并为疾病识别、风险评估、治疗决策和健康管理提供辅助。因此，理解医学数据的基本类型、处理流程及其与具体任务之间的对应关系，是开展医学人工智能应用的重要基础。


\subsection{医学数据类型与基本特征}

患者从入院检查、住院治疗到出院随访的过程中，可能形成电子病历、检验结果、医学影像、连续监护波形和组学检测结果等多种资料。医生需要结合这些信息判断疾病状态、观察治疗反应并评估恢复进程，人工智能则需要将不同来源的数据转换为能够识别、比较和计算的形式，才能从中提取与诊断、治疗和预后相关的特征。

按照数据来源和表现形式，医学数据通常可分为电子病历数据、医学影像数据、生理信号数据和组学数据等类型。电子病历同时包含诊断编码、检验数值等结构化信息，以及主诉、病史和诊疗记录等自然语言文本；医学影像以二维图像、三维体数据或视频形式呈现组织结构；生理信号具有连续采集和随时间变化的特点；组学数据通常包含大量基因、蛋白质或代谢物等高维特征，可用于疾病分型和生物标志物筛选。不过，组学检测容易受到样本量、检测平台、批次差异和统计校正方法影响，相关结果进入临床应用前需要经过独立验证。

如图\ref{fig:ch9_01}所示，不同医学数据在结构形式、时间尺度、信息密度和临床含义方面存在明显差异，对应的人工智能处理方法也各不相同。文本数据常用于信息抽取和临床文本分类，影像数据适合开展分类、检测和分割，连续生理信号可用于状态识别与异常预警，组学数据则可支持生物标志物筛选和疾病分型。多种数据相互补充，共同为疾病识别、病情评估、治疗决策和长期随访提供信息基础。


\par
\noindent\begin{minipage}{\textwidth}
\centering
\BookFigureImage{./images_chapter_07-08/9.1.png}
\captionof{figure}{医学数据类型与基本特征示意图}
\label{fig:ch9_01}
\end{minipage}
\par

\subsubsection*{\textbf{•电子病历数据}}
门诊医生在电脑中写下患者的主诉、现病史和初步诊断时，一份电子病历就开始形成。电子病历数据是现代医疗信息系统中最基础、最常见的一类数据，主要来源于门诊、住院和随访等临床活动。早期电子病历更多是把纸质病历搬到电脑上，主要解决保存和查阅问题；随着医学信息化不断发展，电子病历已逐渐成为整合患者诊疗全过程的重要平台，记录着从入院评估、检验检查、治疗护理到出院随访的完整轨迹。

电子病历中既包含结构化数据，也包含非结构化文本。结构化数据通常以固定字段保存，例如年龄、性别、体温、血压、检验数值、诊断编码和药物剂量等，计算机可以较容易地进行统计和比较。非结构化文本则主要来自医生书写的病程记录、主诉、现病史、出院小结和随访记录等内容，这些文字更接近真实临床表达，包含丰富的医学判断和诊疗线索。例如，一句“活动后胸闷较前加重，夜间不能平卧”，不仅记录了症状，还隐含了疾病进展和心功能状态变化的信息。人工智能要利用这类文本，通常需要借助自然语言处理技术，从中识别症状、疾病、药物、检查、治疗措施和时间关系等关键信息。从内容上看，电子病历通常包括以下几类信息：

\begin{itemize}
\item
  患者基本信息，如姓名、年龄、性别、身高、体重、联系方式等
\item
  主诉、现病史、既往史、个人史和家族史等病史信息
\item
  体格检查结果与临床观察记录
\item
  检验检查结果，如血常规、生化指标、凝血功能等
\item
  诊断、治疗、手术、护理及用药记录
\item
  出院小结、随访记录和康复记录等
\end{itemize}


\indent 电子病历连续记录患者从初次就诊、检查诊断、治疗处置到出院随访的诊疗过程，其中既包含主诉、病史、诊断和治疗计划等文本信息，也包含检验结果、用药记录和生命体征等结构化数据，可为临床决策、医疗管理和人工智能分析提供重要依据。


\subsubsection*{\textbf{•医学影像数据}}
在影像科工作站前，医生每天要面对大量X线片、CT切片、MRI序列、超声图像和病理图像。医学影像数据能够把人体内部结构、器官形态及部分功能状态转化为可观察的图像，是疾病筛查、病灶发现、范围评估和疗效观察中的重要依据。随着成像设备和成像技术不断进步，医学影像的种类持续丰富，数据规模也迅速增长，不同影像在成像原理、空间分辨率和临床用途上存在差异，常见类型主要包括：

\begin{itemize}
\item
  X射线成像：X射线成像是利用X射线穿透人体后形成组织密度差异图像的医学成像技术。骨骼等高密度组织对X射线吸收较多，在图像中通常呈现较亮区域，而含气组织吸收较少，通常呈现较暗区域。该技术具有检查速度快、操作简便、成本较低和设备普及率高等特点，主要用于骨折筛查、胸部疾病初步检查及部分急诊情况的快速评估。
\item
  DSA（digital subtraction angiography，数字减影血管造影）：数字减影血管造影是通过对造影前后的X射线图像进行数字减影，突出显示血管结构和血流状态的成像技术。该技术能够减少骨骼和软组织对血管显示的干扰，较清晰地呈现血管走行、管腔变化及局部血流情况。DSA主要用于脑血管、心血管和外周血管疾病的诊断及介入治疗，可辅助识别血管狭窄、闭塞、畸形和动脉瘤等病变。
\item
  计算机体层成像（CT）：计算机体层成像是利用X射线从多个角度扫描人体，并通过计算机重建获得断层图像的医学成像技术。CT能够减少组织重叠造成的影响，较清楚地显示器官、组织及病变之间的空间关系，并可进一步生成冠状面、矢状面和三维重建图像。该技术具有扫描速度快、空间分辨率较高和结构显示清晰等特点，广泛用于颅脑出血、肺部病变、腹部急症和骨性损伤的检查与评估。
\item
  磁共振成像（MRI）：磁共振成像是利用强磁场和射频脉冲激发人体组织中的氢质子，并根据其产生的磁共振信号重建图像的成像技术。不同组织的含水量、质子密度和弛豫特性存在差异，因此可在MRI图像中呈现不同信号。MRI具有软组织分辨能力强、成像序列丰富和不使用电离辐射等特点，适用于神经系统、肌骨系统、关节、脊柱及部分肿瘤性疾病的检查，尤其有利于观察软组织病变和复杂解剖结构。
\item
  核医学成像：核医学成像是将放射性核素示踪剂引入人体，并通过探测其在组织或器官中的分布和代谢活动形成图像的技术。与主要显示解剖结构的成像方式相比，核医学成像更侧重反映组织代谢、血流灌注和器官功能状态。常见技术包括PET（positron emission tomography，正电子发射断层扫描）和SPECT（single-photon emission computed tomography，单光子发射计算机体层成像），主要用于肿瘤定位与分期、神经系统疾病评估、心肌灌注分析和疗效观察。
\item
  B型超声：B型超声是利用高频声波在不同组织界面产生的反射信号，形成二维灰度图像的医学成像技术。该技术能够实时显示器官形态、组织结构及部分运动状态，并具有无创、便捷、可重复检查和可床旁操作等特点。B型超声广泛用于心脏、腹部、甲状腺、乳腺、妇产科和血管系统检查，可辅助观察器官结构、病灶形态及动态变化。
\item
  内镜成像：内镜成像是将带有光源和成像装置的内镜送入人体自然腔道或手术通道，直接获取腔道和器官表面图像或视频的技术。常见设备包括胃镜、肠镜、支气管镜和腹腔镜，不同设备对应不同的观察部位和诊疗需求。内镜图像与视频能够显示黏膜形态、组织颜色、表面纹理和局部病变，主要用于病灶发现、活检定位、术中观察及治疗过程记录。
\end{itemize}

从数据形式看，医学影像主要表现为图像和视频，具有空间结构信息丰富、视觉特征明显和信息密度高等特点。不同影像技术像是从不同角度观察人体：X射线适合骨性结构和胸部初筛，CT适合断层结构观察和急诊评估，MRI更适合软组织和神经系统检查，超声则具有实时、便捷和无创的优势。


\subsubsection*{\textbf{•生理信号数据}}
监护仪屏幕上跳动的曲线，记录的是人体正在发生的变化。生理信号数据是人体生命活动过程中产生并可被采集、记录和分析的信号信息。与电子病历和影像数据相比，生理信号更强调连续性和动态变化，能够反映机体在短时和长时尺度上的状态波动，常见生理信号数据可以按系统进行划分如下：

\begin{itemize}
\item
  循环系统生理信号，如心电信号、脉搏波、血压变化等
\item
  呼吸系统生理信号，如呼吸频率、呼吸波形、血氧变化等
\item
  神经系统生理信号，如脑电信号、诱发电位等
\item
  运动系统生理信号，如肌电信号、步态信号和动作感知数据等
\end{itemize}

生理信号通常表现为连续波形或时间序列，具有实时性强、变化快和动态特征明显等特点。心电图可以帮助识别心律失常，脑电图可用于监测神经活动，监护设备连续记录的血氧和呼吸参数则有助于危重患者状态评估。与单次检查结果相比，连续信号更像一段医学录像，能够显示状态如何变化，也能帮助发现短暂但重要的异常事件。

近年来，可穿戴设备的发展使长期生理信号监测从院内延伸到院外。智能手表、动态心电监测设备、睡眠监测设备和运动手环等产品，使心率、睡眠、活动量等信息能够在日常生活中持续记录。对于慢病管理、康复随访和异常预警而言，这类数据让人工智能有机会观察到医院之外的健康变化。

\subsubsection*{\textbf{•生命组学数据}}
人体的许多变化不仅表现为可见的症状和检查结果，还会留在基因表达、蛋白质活动和代谢物波动之中，这些分子层面的信息如同一组记录生命状态的微观线索。生命组学数据是从基因、转录、蛋白质和代谢物等层面描述人体生物学状态的数据，能够反映疾病发生、发展及治疗反应过程中产生的分子变化。这类数据通常通过高通量测序和MS（mass spectrometry，质谱分析）等技术获得，具有数据规模大、特征维度高和变量关联复杂等特点。通过分析不同分子的组成及其变化规律，可以为疾病机制研究、生物标志物发现、疾病分型和个体化治疗提供依据。根据研究对象不同，生命组学数据主要包括以下几类：

\begin{itemize}
\item
  \textbf{基因组学数据}：基因组学数据以DNA（deoxyribonucleic acid，脱氧核糖核酸）序列及其变异信息为主要内容，用于描述个体遗传物质的组成、结构和差异。通过分析单核苷酸变异、插入缺失和拷贝数变化等信息，可以为遗传性疾病识别、肿瘤突变检测、疾病易感性评估和个体化诊疗提供依据。

\item
  \textbf{转录组学数据}：转录组学数据记录特定细胞、组织或机体在一定条件下核糖核酸的种类和表达水平，能够反映基因在不同生理状态或病理状态下的表达活动。通过比较不同组织、疾病阶段或干预条件下的表达差异，可以分析基因功能变化，并为疾病分型、作用机制研究和潜在生物标志物筛选提供信息。

\item
  \textbf{蛋白质组学数据}：蛋白质组学数据描述细胞、组织或机体中蛋白质的种类、丰度、结构修饰及相互作用，能够反映基因表达产物参与生命活动的实际状态。通过分析蛋白质表达及功能变化，可以揭示疾病相关生物过程，为疾病机制研究、生物标志物发现和治疗靶点筛选提供依据。

\item
  \textbf{代谢组学数据}：代谢组学数据记录机体在特定生理或病理状态下小分子代谢物的种类、含量及变化规律，能够较直接地反映细胞代谢活动和机体功能状态。通过分析代谢物及其相关代谢通路的变化，可以为疾病诊断、药物作用评价、疗效监测和个体健康状态评估提供参考。
\end{itemize}

与主要记录症状、体征和检查结果的传统临床数据相比，生命组学数据更关注机体在基因、转录、蛋白质和代谢物层面的变化，能够为解释疾病发生机制和个体生物学差异提供分子依据。不同类型的组学数据反映的信息各有侧重，其中，基因组学数据用于识别遗传变异和肿瘤突变，转录组学数据用于分析基因表达变化，蛋白质组学数据用于观察蛋白质表达及功能状态，代谢组学数据则用于反映机体代谢活动。将这些信息综合分析，可以支持疾病机制研究、疾病分型、生物标志物与治疗靶点发现，并为个体化干预提供参考。

\subsubsection*{\textbf{•数据特征}}

医学数据产生于患者就诊、检查、治疗、监测、康复和随访的全过程，其内容既包括数值、文本和图像，也包括波形、视频及生命组学信息。不同数据在采集方式、组织结构和医学含义上存在明显差异，因此需要采用不同的方法进行处理和解释。总体来看，医学数据主要具有多源异构性、专业与情境依赖性、动态时序性以及安全与伦理敏感性，这些特征共同影响医学人工智能的数据处理方式、模型训练过程和结果应用范围。

医学数据首先具有多源异构性。多源是指同一患者的数据可能来自门诊问诊、实验室检查、影像检查、护理记录、监护设备和院外随访等不同环节；异构则是指这些数据分别表现为结构化数值、自然语言文本、医学图像和连续时间序列。以一名突发胸痛患者为例，急诊医生需要同时查看病历中的疼痛部位和持续时间，观察心电图上的波形变化，等待血液检查中的肌钙蛋白结果，并结合冠状动脉影像判断血管是否狭窄。屏幕上的曲线、检验单上的数值和影像中的血管结构，虽然形式不同，却共同指向患者当前的心脏状态。

多源数据之间还具有相互补充和验证的作用，单一指标往往只能反映疾病的某个侧面。以肺部感染为例，患者出现发热和咳嗽，说明呼吸系统可能存在异常；血液检查显示白细胞计数和CRP（C-reactive protein，C反应蛋白）升高，提示体内可能发生炎症反应；胸部CT中出现斑片状阴影，则进一步提供病变的位置和范围。症状、检验指标和影像结果如同几块相互关联的拼图，只有将它们组合起来，才能形成较完整的判断，这也说明医学人工智能需要具备多源数据融合能力。

医学数据还具有较强的专业性和情境依赖性，同一项指标在不同患者和不同诊疗阶段中可能具有不同意义。白细胞计数升高并不总是等同于感染，若患者同时出现高热、咳嗽和肺部阴影，感染的可能性较高；若这一变化发生在剧烈运动后、手术后或使用糖皮质激素期间，则可能与应激反应或药物作用有关。血糖升高同样需要结合具体情境判断，糖尿病患者的持续升高可能提示控制不佳，而严重创伤患者短时间内出现高血糖，也可能是机体应激反应的结果。

医学数据具有明显的动态时序性，患者的健康状态和疾病过程会随时间持续变化。许多临床问题关注的不是某一次检查结果，而是指标在一段时间内的变化趋势。以心肌损伤评估为例，患者刚到急诊时肌钙蛋白可能仍在正常范围，但数小时后的复查结果明显升高，便可能提示心肌损伤正在发生。连续血糖监测也能够把原本分散的测量结果连接成完整曲线，使餐后升高、夜间低血糖和全天波动规律清楚呈现；肿瘤患者治疗前后的多次CT检查，则可以通过病灶大小变化辅助评价治疗效果。

医学数据同时具有较高的安全与伦理敏感性。患者姓名、身份证号码、疾病诊断、遗传检测结果、用药记录和预后情况均涉及个人隐私，一旦在存储、传输或共享过程中管理不当，可能造成信息泄露并损害患者权益。医学影像文件看似只包含一张图像，实际上还可能保存患者姓名、检查编号和检查时间等信息，因此研究人员在使用这类数据训练模型前，需要删除能够识别患者身份的内容，并完成伦理审查和数据授权。医院还应根据岗位职责设置访问权限，使不同人员只能查看工作所需的数据，从而保证医学数据在合法、安全和可追溯的条件下使用。


\subsubsection*{\textbf{•小结}}

本节介绍了医学数据的主要类型及其基本特征。电子病历、医学影像、生理信号和生命组学数据从不同层面记录人体健康状态，共同构成医学人工智能应用所依赖的数据体系。同时，医学数据具有多源异构、专业性强、动态变化明显和敏感性高等特点。理解这些类型与特征，有助于进一步认识数据处理流程、医疗任务需求和人工智能医学应用方式。

\subsection{数据处理流程与标准规范}
一份医学数据从诊疗现场产生，到真正成为可以分析和建模的信息，还需要经历一系列整理与转换。原始记录往往分散在不同医疗系统中，它们可能来自不同设备，采用不同名称、单位和文件格式，还可能夹杂缺失记录、重复内容或采集噪声等问题。只有经过规范处理，才能准确还原患者的健康状态。

医学数据处理主要包括采集、清洗、标注和标准化四个环节，其中，数据采集负责完整保存诊疗信息，数据清洗用于识别并处理缺失、重复和异常内容，数据标注将医学专家的判断转化为模型能够学习的标签，数据标准化则统一术语、格式、单位和编码规则。数据处理解决原始记录如何转化为可用信息的问题，标准规范保证不同来源的数据能够被准确理解、比较和共享。如图\ref{fig:ch9_02}所示，这些环节前后衔接，共同保障医学数据的完整性、一致性和可用性。

\par
\noindent\begin{minipage}{\textwidth}
\centering
\BookFigureImage{./images_chapter_07-08/9.2.png}
\captionof{figure}{医学数据处理流程示意图}
\label{fig:ch9_02}
\end{minipage}
\par

\subsubsection*{\textbf{•数据采集}}

急诊室的门刚刚打开，一名胸痛患者便被送入抢救区。医生询问疼痛部位和持续时间，心电图机开始记录波形，监护仪持续显示心率、血压和血氧，血液样本也被送往检验科。短短几分钟内，文字、数值和波形等信息接连产生，这些记录需要及时进入相应系统，才能形成完整的诊疗数据。

医学数据采集主要包括患者基本信息、诊疗记录、检查检验结果、连续监测数据以及随访和院外健康信息。不同数据采用的采集方式并不相同，病历和用药信息通常由医务人员录入，检验结果、医学影像和监护信号多由专业设备自动生成。以高血压患者为例，门诊病历记录症状和用药情况，家庭血压计持续记录血压变化，随访记录则反映治疗后的长期效果。

采集过程中需要保留患者标识、采集时间、设备来源和检查参数。对于胸痛患者，症状出现、心电图检查、抽血和用药的时间顺序需要清楚记录，便于判断指标变化发生在治疗前还是治疗后。医学影像还应保存扫描层厚、成像序列和设备型号，例如两次肺部CT使用不同层厚时，病灶显示差异可能来自采集条件，而不一定代表病情变化。

连续监测数据还需要记录采样频率、监测时长和设备状态。例如，可穿戴设备佩戴过松时，心率数据可能出现异常跳变；监护电极接触不良时，心电波形也可能突然中断。这些情况应及时标记，避免将设备问题误认为真实生理变化。

涉及患者身份、影像资料、基因检测和院外监测信息时，还需要明确授权范围、访问权限和保存方式。例如，用于教学、科研或模型训练的数据，应删除能够直接识别患者身份的信息，并保留必要的数据来源和处理记录，为后续清洗、标注和模型训练提供可靠基础。


\subsubsection*{\textbf{•数据清洗}}

医院信息系统中，同一名患者可能在急诊、住院、检验和影像系统中留下多份记录。屏幕上看似整齐的数据，实际可能混有重复报告、缺失字段、错误单位和异常数值。若这些问题未经处理便进入分析，模型可能重复计算同一患者，也可能把记录错误当作疾病特征，因此原始数据需要先经过清洗。

数据清洗是识别并处理重复、缺失、异常、错误和格式不一致等问题的过程。医学数据清洗不能只按数值范围机械删除记录，还需要结合采集时间、设备状态和临床背景进行判断。某项指标明显偏高，可能来自录入错误，也可能反映患者病情突然恶化，两种情况不能采用相同的处理方式。

重复记录通常通过患者编号、就诊时间、检查项目、样本编号和报告时间进行核对，以血常规检查为例，同一份报告可能因系统同步和人工补录出现两条相同记录，若全部保留，该患者的数据便会被重复计算。不过，同一患者在不同时间完成的两次检查并不属于重复数据，它们可能反映病情变化，需要按照采集时间分别保存。

缺失数据需要先判断缺失原因，再选择处理方法。身高、体重等一般变量可根据任务需要采用合理的填补方法，关键诊断、主要结局和重要时间点缺失时，则应单独标记或排除相关样本。病历中未填写药物过敏史，并不表示患者不存在过敏，而可能只是信息没有录入；若直接填为“无过敏”，便会把未知状态错误地改成确定结论。

异常值的处理更依赖医学背景，例如，血钾明显升高可能由样本溶血、单位错误或录入错误造成，也可能提示严重电解质紊乱；血氧饱和度突然下降，既可能是监护探头脱落，也可能反映真实的呼吸功能恶化。清洗时需要结合仪器范围、样本状态、相邻时间点和临床记录进行核查，确认异常值的来源后再决定修正、删除或保留。

格式统一主要处理字段名称、计量单位、时间格式、编码方式和文本表达之间的差异。以血糖为例，部分系统采用mmol/L，另一些系统采用mg/dL，若不先换算便直接比较，模型会得到错误结果；“WBC（white blood cell count，白细胞计数）”和“白细胞”也需要映射到同一字段，避免相同指标被识别为多个变量。

不同类型的数据需要采用不同的清洗方法，病历文本主要处理重复模板、错别字和不规范缩写，检验数据重点核对单位、参考范围和样本状态，医学影像需要检查文件完整性、图像方向和严重伪影，生理信号则应识别噪声、信号中断和时间戳错误。以心电数据为例，电极接触不良产生的高幅噪声应被标记，但短暂出现的异常心律不能因偏离正常波形而被直接删除。

数据清洗过程中还应记录删除、填补、修正和单位转换等操作，使后续分析能够追溯数据变化。清洗规则发生调整时，也应保留不同版本及其处理结果，便于比较模型表现并复核关键样本。经过清洗的数据仍需进入标注和标准化环节，逐步形成适合医学分析和模型训练的数据集。

\subsubsection*{\textbf{•数据标注}}

影像科阅片室里，医生在一组肺部CT图像中发现可疑阴影，随后圈出结节位置，并沿着可见边缘勾画轮廓。经过这一操作，原本只由灰度像素组成的图像被赋予了明确的医学含义，模型也由此获得了可以学习的参考信息。数据标注是依据专业知识，对数据中的类别、位置、边界或关键事件进行标记，将医生、技师或研究人员的判断转化为模型训练和评价所需的标签。

根据标记内容的不同，医学数据标注通常采用类别标注、空间标注和时间标注三种形式，其中，类别标注用于说明样本属于“正常”“异常”“良性”或“恶性”等类别，空间标注用于确定病灶位置、器官轮廓和细胞核区域，时间标注则用于指出异常事件发生的时刻或持续区间。以肺结节分析为例，“有结节”或“无结节”的标签可用于图像分类，结节位置框可用于目标检测，沿病灶边缘勾画的完整轮廓则可以支持图像分割和体积测量。

医学影像、病历文本和生理信号在内容结构上存在差异，标注对象也需要随之调整。医学影像通常标记病灶位置、边界和解剖结构，病历文本侧重识别症状、诊断、药物和治疗措施，生理信号则需要标出异常波形及其起止时间。例如，在“胸闷较前缓解，继续口服阿司匹林”这句话中，“胸闷缓解”表示症状变化，“阿司匹林”属于药物信息；在长程心电记录中，除了判断是否出现异常节律，还需要标明异常片段的开始和结束位置，使模型同时学习波形特征和时间变化。

为了减少不同标注人员之间的判断差异，正式标注前需要明确标注对象、类别定义、边界处理方法和不确定样本的处理方式。肺结节边界模糊时，应提前规定按照可见边缘还是估计范围进行勾画；心电信号中出现异常波形时，也需要结合相邻波形和设备状态，区分真实心律异常与电极接触不良造成的干扰。统一的标注规则能够减少个人经验带来的偏差，并为后续复核提供共同依据。

同一数据由不同人员标注时，结果之间可能存在差异，因此实际工作中常采用双人独立标注、专家复核、抽样质检和一致性评价等方式控制质量。若两名医生勾画出的肿瘤范围差异明显，模型便难以形成稳定的边界判断，此时应由经验更丰富的专家结合原始影像和临床资料进行复核。对于肿瘤分割、病理分级和高风险疾病筛查等任务，关键样本通常需要经过多轮确认，以减少错误标签对模型训练的影响。

标注内容的精细程度应根据模型任务和实际需求进行选择，类别标签制作较快，适合完成整体分类；位置框能够指出病灶所在区域，可用于目标检测；精确轮廓则可以支持面积、体积和形态分析，但需要投入更多专业时间。数字病理图像中的细胞核标注、肿瘤边界逐像素勾画以及长程脑电中的异常片段标记，都属于工作量较大的精细标注，因此不宜脱离研究目标盲目增加标注复杂度。

部分病例本身处于早期、边界或难以确定的状态，标注时可以设置“不确定”“待复核”或“低置信度”等标签，也可以保留多名专家的独立意见，避免强行归入某个确定类别。标注结果还应记录标注人员、标注时间、规则版本和复核情况，便于后续追溯标签来源和处理过程。经过规范标注的数据能够用于模型训练、性能评价和结果分析，并使模型输出具有明确的医学参照。


\subsubsection*{\textbf{•数据标准化}}

两家医院调取同一位患者的检验记录时，屏幕上可能出现“WBC”“白细胞计数”和“白细胞”三种名称，虽然医生能够判断它们表示同一项检查，计算机却可能将其识别为三个不同变量。数据在跨医院传输、合并或用于模型训练前，需要按照共同规则统一表达方式，使相同医学概念能够准确对应。

数据标准化主要对清洗后的数据进行术语映射、编码转换和结构整理，使来自不同系统的记录具有一致的医学含义。它与数据清洗的侧重点不同，清洗主要处理错误、缺失和重复内容，标准化则解决同一概念如何采用统一方式表示的问题。例如，“心梗”“心肌梗死”和“急性心肌梗死”需要根据具体语境映射到相应的规范名称，避免模型因表述差异而遗漏相关病例。

医学术语与编码体系能够为不同系统提供共同的识别规则。疾病诊断可以按照ICD（International Classification of Diseases，国际疾病分类）进行编码，医学影像通常采用DICOM（Digital Imaging and Communications in Medicine，医学数字成像与通信）格式保存和传输。以急性心肌梗死为例，不同医生可能采用中文全称、简称或英文缩写进行记录，统一编码后，检索和统计过程便可以将这些记录关联到同一类疾病。

医学数据还需要采用一致的时间表达和数据结构，尤其是涉及诊疗顺序和连续监测的任务。入院、检查、用药和手术等时间应采用统一格式，并保留能够满足分析要求的时间粒度。例如，研究降压药物的起效过程时，只有日期而缺少具体用药时刻，便难以将血压变化与治疗时间准确对应；对于心电或血糖等连续数据，还需要统一时间戳和采样间隔，避免不同设备的记录发生错位。

不同类型的数据在标准化过程中需要保留与其医学含义相关的信息。病历文本应对应规范术语，检验数据应关联统一项目名称和参考范围，医学影像应保留扫描参数和序列信息，生命组学数据则需要统一样本编号和分子命名。两次CT检查的扫描层厚不同，图像差异可能来自采集条件，因此不能在统一格式时删除这些参数。

标准化并不是将所有数据强行改成完全相同的形式，而是在统一表达的同时保留必要的原始值、来源信息和转换规则。不同实验室的参考范围可能因检测方法和试剂不同而存在差异，这类信息应与标准化结果一并保存，便于后续解释模型结果和追溯数据来源。经过标准化的数据能够更顺利地进入多源融合、模型训练和跨机构验证等环节。


\subsubsection*{\textbf{•小结}}

本节介绍了医学数据进入人工智能分析前的基本处理流程，包括数据采集、清洗、标注和标准化等环节。数据采集决定原始信息能否被完整保存，数据清洗用于处理错误、缺失和异常内容，数据标注赋予数据明确的医学含义，数据标准化则使不同来源的数据能够按照统一规则记录和使用。上述环节相互衔接，共同构成医学数据进入分析、研究和实际应用的基础。


\subsection{多源数据融合与表示学习}

一名疑似脑卒中的患者被送入急诊后，医生会迅速询问发病时间、肢体活动情况和既往病史，同时结合头颅CT判断是否存在脑出血，并通过磁共振成像观察缺血区域。血压、血糖和凝血功能等指标也会陆续返回，为治疗方案选择提供补充依据。患者的病情并不完整地存在于某一张图像或某一项检验结果中，而是分散在病历文本、医学影像和检验数据之中，只有将这些信息放在同一诊疗背景下综合分析，才能形成较完整的患者状态描述。

多源数据融合是将不同系统、设备和诊疗环节产生的数据进行关联与整合，使模型能够同时利用多方面信息；表示学习则将文本、图像、波形和分子变量等原始数据转化为模型可以计算和比较的特征。例如，模型可以从病历文本中提取症状和病史特征，从磁共振图像中提取缺血区域特征，再结合血压和凝血指标进行风险判断。由于这些数据的记录时间、表现形式和数值范围并不相同，实际处理通常包括三个相互衔接的环节：首先按照患者身份和诊疗时间完成数据对齐，其次分别提取不同类型数据的特征表示，最后将这些特征进行融合并用于后续分析，如图\ref{fig:ch9_03}所示。

\par
\noindent\begin{minipage}{\textwidth}
\centering
\BookFigureImage{./images_chapter_07-08/9.3.png}
\captionof{figure}{多源数据融合与表示学习示意图}
\label{fig:ch9_03}
\end{minipage}
\par

\subsubsection*{\textbf{•多源数据类型}}
一名患者从入院检查到治疗随访，会在不同系统和诊疗环节中留下相互关联的记录。这些数据可能产生于同一次就诊，也可能分布在不同时间阶段，其记录频率和信息粒度并不相同。例如，病历记录通常围绕一次诊疗过程展开，医学影像反映某一检查时刻的组织结构，监护设备则可能以秒为单位连续记录生命体征。多源数据分析需要将这些分散的信息对应到同一患者和同一临床过程，避免将不同时间或不同诊疗阶段的数据错误组合。

不同来源的数据往往从不同角度描述同一疾病状态，并在分析中形成互补。以肺部感染评估为例，病历文本记录发热、咳嗽和咳痰等症状，血常规和C反应蛋白反映炎症状态，胸部CT显示病变位置和累及范围，血氧变化则提示呼吸功能是否受到影响。单独分析影像时，模型可能难以区分感染、肿瘤或其他肺部异常；仅使用检验指标时，又无法获得病灶位置和形态信息。将这些数据结合起来，可以同时利用症状、炎症指标、影像表现和生理变化，更完整地描述患者状态。

\subsubsection*{\textbf{•数据对齐处理}}

一名患者在急诊完成头颅CT后，检验结果陆续返回，随后接受药物治疗并转入住院病房。数小时内产生的影像、检验指标、生命体征和医嘱记录分散在不同系统中，如果只按患者姓名或检查日期进行合并，治疗后的数据可能被误配到治疗前，甚至与该患者另一次住院记录混在一起。数据对齐就是根据患者、就诊事件和时间关系，将这些记录重新对应到同一诊疗过程。

实际处理中，需要先通过患者唯一标识关联不同系统中的记录，再结合就诊编号、入院时间和科室信息区分不同诊疗事件。同一名患者可能在一年内多次住院，每次住院又包含多次检查和用药记录，因此患者身份相同并不表示数据属于同一次就诊。以糖尿病患者为例，门诊复查血糖与住院期间的连续血糖记录应分别归入相应诊疗阶段，避免将长期随访数据错误地并入急性治疗过程。

不同数据的产生频率差异较大，时间对齐时通常需要确定临床事件和观察窗口。监护设备可以每秒记录心率和血氧，检验指标可能数小时检测一次，医学影像则只在特定时刻采集。分析降压药物效果时，可以把首次用药作为时间基准，将用药前后的血压记录划入相应时间窗口；分析脑卒中患者时，则需要将发病、影像检查、溶栓治疗和复查结果按照先后顺序组织，避免忽略治疗时机对结果的影响。

当连续信号与低频检查结果需要联合分析时，还要处理时间粒度不同的问题。例如，监护仪在一小时内产生数千个心率数据点，而同一时段可能只有一次血液检验结果，此时可根据任务提取该时间窗口内的均值、最大值、波动范围或异常持续时间，再与检验结果建立对应关系。过度压缩会掩盖短暂异常，时间窗口过窄又可能遗漏相关信息，因此窗口长度应结合疾病过程和分析目的确定。

对于时间戳缺失、设备时钟偏移或诊疗阶段不明确的记录，不宜强行建立对应关系。系统时间比实际采集时间晚数分钟时，应根据设备日志或操作记录进行校正；无法确认属于治疗前还是治疗后的数据，则需要单独标记或排除。完成对齐后，每组数据应能够明确对应患者、就诊事件、临床阶段和时间窗口，使后续特征提取与数据融合建立在正确的诊疗顺序上。


\subsubsection*{\textbf{•表示学习方法}}

医生打开一名胸痛患者的诊疗资料时，屏幕上既有病历文字和检验数值，也有冠状动脉CT图像与连续心电波形。医生可以依据专业知识理解这些信息之间的联系，模型则需要先把文字、图像、波形和分子变量转化为可计算的特征。表示学习就是从原始数据中提取具有分析价值的特征，使不同形式的医学信息能够进入模型进行比较、分类和预测。

医学数据的特征表示通常包括人工特征、降维特征、深度学习特征和多源联合特征。人工特征主要由研究人员根据医学知识选取或计算，例如从肺结节图像中提取大小、密度、边缘和纹理，从心电波形中计算心率、节律和波形间期，从结构化病历中选取年龄、血压和检验指标。这些特征含义清楚，便于医生理解，但面对复杂图像、长时间信号和高维分子数据时，人工设计往往难以覆盖全部信息。

当原始变量数量较多且相关性较强时，可以通过降维方法提取少量综合特征。PCA（principal component analysis，主成分分析）能够将多个相关变量转换为少数主成分，从而减少信息冗余和模型计算量。以基因表达数据为例，一份样本可能包含上万个基因变量，直接建模容易受到噪声和样本量不足的影响，通过降维或特征筛选可以保留主要变化，并降低模型过拟合的风险。

深度学习方法能够直接从原始数据中逐层学习特征，其中，CNN（convolutional neural network，卷积神经网络）常用于医学影像分析，可以从像素中识别边缘、纹理和病灶结构；RNN（recurrent neural network，循环神经网络）及其相关结构可用于心电、脑电和连续监测数据，提取随时间变化的信号规律；Transformer模型则适合处理病历文本和较长的序列信息。例如，在胸部CT分析中，网络可以从早期的局部纹理逐步形成结节形态特征，而不需要研究人员逐项设定病灶边缘和密度指标。

多源数据联合分析时，模型通常先分别提取各类数据的特征，再将它们组合成患者的综合表示。以肿瘤治疗反应评估为例，模型可以从影像中提取病灶大小和形态变化，从病历中提取治疗方案和症状记录，从基因检测中提取突变信息，再结合检验指标分析患者的治疗反应。不同数据的特征维度和信息质量可能存在差异，因此联合表示时还需要避免某一类数据因维度过高或噪声过大而主导模型结果。

特征表示需要保留与医学任务相关的信息，并尽量减少设备型号、记录习惯和偶然噪声的干扰。肺部影像特征应反映病灶结构，而不是只识别某家医院的扫描参数；病历文本特征应关注症状、诊断和治疗关系，而不是依赖特定医生的书写格式；心电特征也应区分真实节律变化与设备噪声。只有当特征能够稳定反映疾病状态和诊疗过程时，表示学习的结果才适合用于后续数据融合和模型训练。


\subsubsection*{\textbf{•数据融合方式}}

肿瘤多学科会诊时，影像科医生展示病灶大小和位置，病理科医生说明组织类型，检验结果呈现患者当前状态，基因检测则补充突变信息。人工智能模型也可以在不同阶段整合这些资料，根据融合发生的位置，常用方法包括早期融合、中期融合和晚期融合。

早期融合是在模型训练前，将不同来源的数据整理成统一的特征表。以肺癌复发风险预测为例，可以把患者年龄、吸烟史、肿瘤分期、血液检验结果和影像组学特征放入同一张表中，再输入模型进行分析。这种方式适合数值形式接近、患者记录较完整的数据，但如果某些患者缺少影像或检验结果，合并后的特征表就容易出现大量空缺。

中期融合先分别处理不同形式的数据，再在模型内部组合各自的特征。以脑卒中评估为例，图像分支从头颅CT中提取出血或缺血特征，文本分支从病历中提取发病时间和症状信息，数值分支处理血压、血糖和凝血指标，随后将这些特征合并用于判断病情。这样的处理能够保留图像、文本和数值数据各自的特点，适合多种医学信息共同参与的任务。

晚期融合则保留多个模型的独立判断，再对输出结果进行整合。例如，肺结节影像模型根据结节形态给出恶性概率，临床风险模型结合年龄、吸烟史和家族史生成风险评分，最终系统可以通过加权平均、投票或预设规则形成综合结果。当影像模型提示高风险而临床模型评分较低时，系统还可以保留两种结论，供医生进一步复核，而不是简单用一个结果覆盖另一个结果。

融合方式应根据数据完整性、样本规模和任务特点进行选择。结构化变量较多且缺失较少时，可以将年龄、检验指标和影像组学特征整理为统一输入，采用早期融合完成分析；当影像、病历文本和生理信号共同参与疾病判断时，中期融合能够分别提取各类数据特征，再在模型内部完成组合；医院已经部署多个独立模型时，则可采用晚期融合整合各模型的预测结果。以小样本肿瘤研究为例，结构复杂的中期融合模型可能因训练数据不足而出现过拟合，此时使用较简单的早期融合，往往更容易获得稳定结果。



\subsubsection*{\textbf{•应用局限}}

一名患者接受肿瘤评估时，影像资料和检验结果可能较为完整，病理报告却尚未返回，基因检测也可能因费用或样本条件无法完成。模型训练时若使用了影像、病理、检验和基因数据，实际应用中缺少其中任何一类信息，都可能使模型无法运行或导致预测结果发生明显变化。因此，多源数据融合需要考虑临床数据并不总是同时产生，也不能保证每位患者都具有完整记录。

不同来源数据的质量差异也会影响融合结果。以连续血氧监测为例，探头松动可能造成数值突然下降，而胸部CT和病历记录并未显示患者呼吸状态恶化；若模型未识别这种设备干扰，错误的血氧数据就可能提高风险评分。影像模糊、病历记录不完整或诊断标签错误也会产生类似影响，使质量较差的数据干扰其他可靠信息。

多源融合模型同时使用文本、影像、信号和结构化指标后，预测依据往往比单一数据模型更难追踪。例如，系统将某名肺结节患者判定为高风险时，医生需要了解这一判断主要来自结节边缘特征、吸烟史还是检验指标异常。若模型只给出风险分数，却不能显示不同数据对结果的影响，医生便难以判断结论是否符合患者的实际情况。

不同医院的设备、检查流程和患者构成还可能存在差异，使模型在新环境中的表现下降。某模型在高分辨率CT数据上完成训练后，应用于扫描层厚较大的影像时，病灶特征可能无法被同样识别。因此，多源融合模型在进入临床应用前，需要检查各类数据的完整程度、质量差异和来源变化，并评估模型在信息缺失及跨机构场景下的稳定性。


\subsubsection*{\textbf{•小结}}
本节介绍了多源数据融合与表示学习的基本内容。医学数据来自电子病历、医学影像、生理信号、生命组学和院外监测等多个来源，不同数据从不同层面记录患者状态。多源数据融合通过数据对齐、特征提取和信息整合，使模型能够综合利用多类医学线索；表示学习则把文本、图像、波形和分子变量等不同形式的数据转化为可计算的特征表达。多源融合有助于疾病识别、风险预测、疗效评价和健康管理，但其效果依赖数据质量、时间对应关系和临床解释，不能脱离医学语境单独使用。

\subsection{医疗人工智能的典型任务}

胸部CT检查完成后，医生需要在连续切片中寻找可疑阴影，确定结节的位置和范围，并结合其大小、边缘和密度判断良恶性。若患者同时存在高龄、吸烟史或慢性肺病，还需进一步评估结节进展和疾病发生风险，并据此安排复查、活检或治疗。随着诊疗过程不断深入，医学分析的内容也由异常发现延伸到疾病判断、风险评估和方案选择。

医疗人工智能通常承担病灶识别、辅助诊断、风险预测和治疗决策支持等任务。病灶识别用于发现并定位影像、信号或病理切片中的异常区域；辅助诊断结合患者表现和检查结果判断疾病类型或严重程度；风险预测利用既往记录和当前指标估计疾病发生、进展或复发的可能性；治疗决策支持则根据患者特征和诊疗资料，为检查选择、治疗方案和随访安排提供参考。以肺结节患者为例，模型可以先标出可疑结节，再判断其恶性风险，并结合患者的临床信息提出复查时间建议。图\ref{fig:ch9_04}展示了这四类任务在医学数据分析和诊疗支持中的主要联系。


\par
\noindent\begin{minipage}{\textwidth}
\centering
\BookFigureImage{./images_chapter_07-08/9.4.png}
\captionof{figure}{医疗人工智能典型任务框架示意图}
\label{fig:ch9_04}
\end{minipage}
\par

\subsubsection*{\textbf{•病灶识别}}

影像科医生逐层浏览胸部CT时，一枚直径只有数毫米的肺结节可能隐藏在血管和正常组织之间；眼科医生检查眼底图像时，也需要从密集血管中辨认细小的出血点和微血管瘤。人工智能病灶识别通过分析医学影像或病理图像中的结构、纹理和形态变化，发现可疑异常并提示其位置，为医生缩小阅片范围和复核重点。

根据输出内容的不同，病灶识别可以形成图像分类、目标检测和图像分割等任务。图像分类用于判断整幅图像是否存在异常，例如判断胸部X线图像中是否出现肺炎征象；目标检测使用边界框标出病灶的大致位置，可用于定位肺结节、骨折和内镜息肉；图像分割则进一步勾画病灶或器官的精确边界，例如分离脑肿瘤区域并计算其体积。任务越精细，标注和训练所需的专业投入也越大。

早期方法通常根据医学经验提取病灶大小、灰度、边缘、纹理和形状等特征，再由分类模型判断区域是否异常。以乳腺X线图像为例，钙化点的形态、数量和分布可以作为识别依据；肺结节的边缘毛刺、密度和圆整程度也能帮助区分不同风险。人工特征的医学含义较清楚，但面对边界模糊、形态多变或背景复杂的病灶时，有限的预设特征往往难以覆盖全部变化。

深度学习方法可以直接从图像中学习与病灶有关的视觉特征，使模型逐步识别局部边缘、组织纹理和整体形态。例如，在肺结节检测中，模型可以先从CT切片中寻找与周围组织不同的局部区域，再结合相邻切片判断其是否形成连续结构；在数字病理图像中，模型则需要从大范围组织背景中识别细胞核异型、排列紊乱和染色异常等局部变化。对于小病灶和低对比病灶，分析不同尺度的图像区域并加强对可疑位置的关注，有助于减少正常组织和背景噪声的干扰。

病灶识别模型的建立通常需要经过图像筛选、病灶标注、预处理、模型训练和结果展示等环节。医生在图像上标出病灶类别、位置或边界后，模型据此学习病变特征；图像进入训练前，还可能进行尺寸调整、灰度归一化和伪影处理。模型完成分析后，可以输出边界框、分割轮廓、热力图或风险概率，医生再返回原始图像核对提示区域。例如，肺结节筛查系统标出可疑位置后，医生仍需结合结节密度、边缘和既往影像判断其临床意义。

病灶识别适用于图像数量大、病灶较小或阅片任务重复度较高的场景。胸部CT分析可以辅助发现肺结节，眼底筛查可以提示出血、渗出和微血管瘤，数字病理分析可以定位肿瘤区域或异常细胞聚集区，内镜图像分析则可标出息肉、溃疡和可疑黏膜病变。这些应用主要帮助医生提高筛查效率和减少遗漏，不能代替病史分析、实验室检查和专业诊断。

不同医院的设备型号、扫描参数、图像质量和标注规则存在差异，模型在训练数据上表现良好后，仍可能在新医院出现漏检或误检。以肺结节检测为例，较厚的扫描层面可能掩盖微小结节，炎症阴影和血管截面也可能被误认为病灶；罕见病灶数量不足时，模型还容易偏向常见类型。因此，病灶识别模型需要通过跨机构数据进行验证，并结合类别平衡、结果解释和医生复核，才能用于实际筛查和辅助诊断。


\subsubsection*{\textbf{•辅助诊断}}

胸部CT中出现一处边缘模糊的肺部阴影后，医生还需要结合患者是否发热、炎症指标是否升高、既往影像是否出现变化以及是否具有长期吸烟史，判断其更接近感染、良性结节还是恶性病变。辅助诊断是在异常区域被发现后，进一步分析病变性质、疾病类别、分型和严重程度，并将分散的医学线索整理为可供医生参考的判断结果。

辅助诊断的输出形式会随临床任务而变化，可以表现为疾病类别、病变等级、分型标签或患病概率。以糖尿病视网膜病变筛查为例，模型不仅需要识别眼底是否存在出血和渗出，还要根据病灶数量、位置和累及范围判断病变程度；在心电图分析中，模型则需要区分房颤、室性早搏等不同异常类型，而不是只提示波形存在异常。相比病灶识别，辅助诊断更加关注异常所代表的医学含义。

临床诊断通常需要将患者表现与检查结果放在同一背景下理解，单项数据只能提供部分依据。肺部阴影可能由感染、肿瘤或炎症后改变引起，若患者伴有发热、咳嗽和C反应蛋白升高，感染的可能性会增加；若阴影持续增大，并伴有边缘毛刺和长期吸烟史，则需要进一步考虑恶性风险。模型可以综合影像表现、病历记录和检验指标，输出疾病类别或风险概率，但这些结果仍需要结合患者的完整诊疗资料进行判断。

辅助诊断方法的选择与数据形式和任务复杂程度有关。年龄、血压、检验指标和病史等结构化变量可以使用传统机器学习方法完成疾病分类或状态判断，医学影像、病理图像、心电信号和病历文本则通常采用深度学习方法提取复杂特征。例如，心血管疾病辅助诊断可以联合胸痛表现、心电变化和肌钙蛋白水平，肿瘤辅助诊断则可结合影像特征、病理类型和基因突变信息。数据种类增加后，模型能够利用的线索更加丰富，但缺失记录和低质量数据也可能干扰最终判断。

模型结果应采用医生能够理解和复核的形式进行展示，除疾病类别或风险概率外，还可以提供关键指标、可疑区域和主要判断依据。例如，肺结节模型给出较高恶性概率时，可以同时显示结节位置、边缘特征和既往体积变化；糖尿病视网膜病变模型给出严重程度等级时，也可以标出影响分级的出血点和渗出区域。只有将结果与原始数据对应起来，医生才能判断模型提示是否符合患者实际情况。

辅助诊断结果受到输入数据质量、训练样本构成和疾病复杂程度的影响，不能直接替代临床诊断。罕见疾病、早期病变和多种疾病同时存在时，模型可能因训练样本不足而出现误判；同一疾病在不同年龄和不同基础疾病患者中，也可能表现出不同特征。部分结论还需要病理检查、动态随访或多学科讨论才能明确，因此辅助诊断主要用于整理医学线索、提示可能风险和提高判断效率，最终结论仍需由医生结合临床资料作出。

\subsubsection*{\textbf{•疾病风险预测}}

一名慢性心力衰竭患者准备出院时，症状已经缓解，血压和血氧也趋于稳定，但医生仍需判断其在未来数周内是否可能再次住院。近期体重是否持续增加、肾功能是否下降、夜间心率是否异常以及过去一年是否多次住院，都可能成为病情再次恶化的早期线索。疾病风险预测通过分析患者当前状态、既往记录和连续变化趋势，估计疾病进展、并发症发生、病情恶化或不良结局出现的可能性。

风险预测需要明确具体的临床结局和预测期限，不同目标对应的数据范围与判断重点并不相同。模型可以预测心衰患者30天内是否再次住院、脑卒中患者1年内是否复发、重症患者未来24小时内是否出现病情恶化，也可以评估肿瘤治疗后的复发风险。例如，预测重症患者短期恶化时，近期血压、血氧和乳酸变化更为重要；评估肿瘤远期复发时，则需要结合病理特征、治疗方式和长期随访记录。

患者状态随时间不断变化，连续记录通常比单次测量包含更多风险信息。以糖尿病管理为例，一次血糖升高可能与进食或短期应激有关，而连续监测中反复出现餐后高峰和夜间低血糖，则提示现有治疗方案可能需要调整。心衰患者若在数日内同时出现体重增加、活动量下降和静息心率升高，这些变化虽然单独看并不突出，联合分析后却可能形成较明确的病情加重信号。

疾病风险预测可以采用传统机器学习方法或时序深度学习方法。对于年龄、病史、检验指标和既往住院次数等结构化数据，可以使用LR（logistic regression，逻辑回归）、SVM（support vector machine，支持向量机）、RF（random forest，随机森林）和GBDT（gradient boosting decision tree，梯度提升决策树）等模型完成风险分层和概率估计。以心衰再入院预测为例，模型可以根据肾功能、血钠、既往住院次数和出院前生命体征，判断患者属于低风险、中风险还是高风险。对于连续心率、血压、血糖和随访记录等时序数据，则可以采用RNN、LSTM（long short-term memory，长短期记忆网络）等方法，提取变量随时间变化的趋势和相互关系。

部分任务还需要同时处理静态信息和动态信息。静态信息包括年龄、性别、基础疾病和既往诊断，动态信息则包括连续生命体征、检验变化和用药调整。以重症患者恶化预警为例，年龄和基础疾病反映患者的基本风险，过去数小时内的血压下降、心率升高和血氧波动则反映当前状态，两类信息结合后可以形成更完整的风险判断。若同时使用影像、病历和连续监测数据，还可以构建多源风险预测模型，但需要注意不同数据之间的时间对应和缺失情况。

预测分析还需要区分观察时间窗和预测时间窗。以30天再入院预测为例，可以使用患者出院前一周的检验结果、生命体征和用药记录作为观察信息，并判断出院后30天内是否再次住院。若把再次住院后产生的检查结果错误地放入观察数据，模型便提前获得了结局发生后的信息，预测结果虽然看似准确，却无法用于真实临床场景。

风险结果通常以发生概率、风险等级、趋势曲线或预警信号呈现，并应同时说明主要影响因素。例如，系统将某名心衰患者列为高风险时，可以提示近期肾功能下降、体重持续增加和既往多次住院等依据。模型输出的概率还需要经过校准，使预测风险与实际事件发生率尽量接近，否则即使分类结果较好，也难以用于随访安排和医疗资源分配。

疾病风险预测提供的是基于已有数据形成的概率判断，不能确定单名患者一定会出现某种结局。治疗调整、用药依从性、突发事件和生活方式变化都会改变实际发展过程，不同医院的患者构成和诊疗方式也可能影响模型表现。因此，预测结果主要用于识别需要重点关注的人群，辅助医生安排复查、随访和早期干预，并根据患者后续状态持续更新。



\subsubsection*{\textbf{•治疗决策支持}}

门诊诊室内，一名同时患有高血压、糖尿病和慢性肾功能不全的患者带着多份检查报告前来复诊。医生调整降压药时，既要观察血压控制情况，也要核对肾功能、血糖变化、既往用药反应和药物禁忌，还需考虑患者能否长期按时服药。治疗决策支持通过整理患者资料、筛选可选方案并提示潜在风险，为医生比较不同治疗策略提供参考。

系统在分析前需要形成与治疗选择相关的患者状态，其中包括主要诊断、合并疾病、症状体征、检验指标、影像结果、肝肾功能、过敏史、既往治疗反应和当前用药等信息。以感染患者为例，抗菌药物选择需要结合感染部位、病原学检查、药敏结果和肾功能；若患者同时存在青霉素过敏或近期使用过同类药物，部分方案便需要排除或调整。

候选治疗方案通常来自临床指南、专家共识、医院诊疗路径、药品说明书和既往病例，并按照适用条件、疗效和风险进行比较。对于肿瘤患者，候选方案可能包括手术、放疗、化疗、靶向治疗和免疫治疗，系统可以根据病理类型、影像分期、基因突变和体能状态，提示各方案的适用条件与主要风险；对于慢性病患者，候选内容则可能涉及药物组合、剂量调整、生活方式干预和复查时间。

治疗决策支持可以通过医学规则和病例数据两类信息形成建议。医学规则适合处理药物过敏、禁忌证、剂量调整和危急值处置等明确要求，例如患者肾功能下降时，系统可以提示某些药物需要减量或停用；KG（knowledge graph，知识图谱）还可以连接疾病、药物、不良反应和检查项目，用于发现潜在冲突。病例数据则用于分析患者特征、治疗方案和临床结局之间的关系，例如比较具有相似病理类型和基因突变的肿瘤患者接受不同治疗后的反应。

对于治疗过程会随患者状态连续调整的场景，系统还可以根据新的监测结果更新建议。重症患者接受液体治疗后，血压、尿量和乳酸水平会持续变化，后续处理需要根据这些指标重新评估；糖尿病患者的胰岛素剂量也可能随着进食、运动和血糖波动调整。RL（reinforcement learning，强化学习）可以研究这类连续决策过程，但模型形成的策略必须经过严格的离线验证和临床评估，不能直接在真实患者中试错。

系统输出通常包括候选方案排序、用药风险提示、剂量调整建议和随访安排，并应同时说明主要依据。以抗菌药物管理为例，系统可以提示推荐药物、适用剂量及需要监测的肝肾功能指标；若存在药物过敏、耐药风险或潜在相互作用，还应明确显示相关原因，而不是只给出单一的“推荐”或“不推荐”结论。

真实患者常存在多病共存、药物可及性不足、依从性较差和个人偏好不同等情况，模型评分较高的方案未必适合直接采用。部分建议还会受到知识更新速度、数据完整性和样本代表性的影响，过多低价值提示也可能造成提醒疲劳。因此，系统应把关键依据、适用条件和风险边界一并呈现，由医生结合患者意愿、医疗条件和治疗目标作出最终选择。


\subsubsection*{\textbf{•小结}}

本节介绍了医疗人工智能的典型任务，包括病灶识别、辅助诊断、疾病风险预测和治疗决策支持。病灶识别侧重发现异常对象，辅助诊断侧重判断疾病性质和类别，风险预测用于估计患者未来状态，治疗决策支持则为干预方案比较与选择提供辅助。不同任务在数据形式、分析目标和结果表达方面存在差异，需要根据具体应用选择合适的机器学习或深度学习方法进行建模。


\subsection{医学数据安全与伦理规范}

医院信息系统中的一份病历，可能同时记录患者姓名、身份证号、联系方式、疾病诊断、检查结果、用药情况和手术记录。影像资料被用于科研分析时，即使删除了姓名，图像编号、检查时间或其他关联信息仍可能暴露患者身份；基因检测数据一旦泄露，还可能影响患者及其家庭成员。医学人工智能需要使用大量临床数据，但数据进入共享、研究和模型训练环节后，访问人员和使用范围随之增加，隐私泄露、越权使用和信息滥用的风险也会相应上升。

医学数据安全是指在数据采集、存储、传输、使用和共享过程中，通过技术与管理措施防止数据被非法访问、泄露、篡改或丢失。电子病历、医学影像、基因检测结果和院外监测记录均包含与个人健康密切相关的信息，因此需要根据数据敏感程度设置访问权限、使用范围和保存期限。例如，科研人员使用临床数据训练模型前，应删除能够直接识别患者身份的信息，并限制数据复制、下载和外传。

医学数据管理还涉及知情同意、公平使用、责任划分和患者权益保护。患者需要了解数据将用于何种目的、由哪些人员使用以及是否会被再次共享，模型开发者则应避免将数据用于超出授权范围的任务。若训练数据主要来自特定地区或特定人群，模型还可能对其他患者产生不公平结果，因此医学人工智能应用既要保护数据本身，也要审查数据使用方式和模型结果可能带来的影响。

\subsubsection*{\textbf{•隐私保护}}

一份用于肺部影像研究的CT数据即使删除了患者姓名，图像中的检查编号、拍摄时间和医院信息仍可能与其他记录发生关联，从而暴露患者身份。隐私保护需要降低数据在使用和共享过程中被重新识别的风险，其中，数据去标识化通常通过删除、替换或模糊姓名、身份证号、电话号码、住址和住院号等信息，使数据能够在限定范围内用于科研、教学和模型训练。

不同数据的识别风险并不相同，处理方式也需要有所区别。普通检验记录可以通过移除直接身份信息降低风险，基因组数据、罕见病资料和具有特殊影像表现的病例则可能因信息组合过于独特而再次指向具体患者。例如，某地区仅有少数罕见病患者时，即使数据中没有姓名，年龄、就诊医院和疾病类型的组合仍可能暴露身份，因此还需要限制访问人员、使用场景和数据共享范围。

\subsubsection*{\textbf{•授权管理}}

临床医生查看患者完整病历是为了完成诊疗，科研人员训练模型时却未必需要接触姓名、联系方式和全部就诊记录，因此不同人员应获得与任务相匹配的数据权限。授权管理通常按照最小必要原则设置访问范围，即只提供完成当前医疗、科研或管理任务所需的数据，并根据岗位区分查看、下载、修改和导出等操作权限。

数据权限还应细化到具体的使用场景和操作范围，例如，研究人员开展影像模型开发时，可以访问去除身份信息后的图像和诊断标签，但不应接触患者联系方式及与研究无关的病历内容；外部合作机构需要使用数据时，还应明确使用目的、保存地点、使用期限和销毁方式。系统可以通过账号认证、分级授权和操作日志记录访问过程，并在出现批量下载、异常登录或越权查询时及时发出警告，减少数据被复制、转交或用于未授权任务的风险。


\subsubsection*{\textbf{•伦理审查}}

研究团队准备利用历史病历训练疾病预测模型时，即使相关数据已经保存在医院系统中，也不能直接将其视为可以自由使用的研究材料。医学伦理委员会需要评估研究目的、数据来源、患者权益、隐私保护措施和潜在风险，并判断研究方案是否符合知情同意和数据使用要求。

知情同意要求患者了解数据的使用目的、可能风险、保护措施和退出方式，并在充分理解后自主作出决定。例如，采集患者基因数据开展长期研究时，应说明数据是否会用于后续项目、是否可能与其他机构共享，以及研究结束后如何保存。对于难以逐一重新取得同意的历史数据，也需要根据风险程度、去标识化情况和研究必要性进行伦理评估，不能仅因数据已经存在或技术上可以调用，便省略审查程序。

\subsubsection*{\textbf{•算法公平}}

某疾病预测模型若主要使用大型医院患者数据进行训练，可能更熟悉检查资料完整、疾病表现典型的病例，而对基层医院中检查项目较少或疾病表现不典型的患者判断不准。算法公平性要求模型在不同年龄、性别、地区、疾病类型和人群来源中尽量避免持续偏向某一群体，使技术应用不会扩大原有的医疗差异。

模型开发时需要检查训练数据是否覆盖主要使用人群，并分别评价不同群体中的准确率、漏诊率和误报率。例如，皮肤病识别模型若缺少不同肤色人群的数据，可能在部分患者中出现更高漏诊率；老年患者和慢性病患者在训练样本中数量不足时，风险预测结果也可能偏低。发现明显差异后，需要通过补充样本、调整训练策略和开展外部验证进行修正，并在应用时说明模型适用范围，避免将某一人群中得到的结果直接推广到所有患者。


\subsubsection*{\textbf{•小结}}
本节介绍了医学数据安全与伦理规范的基本内容。医学数据包含患者身份、健康状况和诊疗过程等敏感信息，在人工智能应用中必须重视隐私保护、授权管理、伦理审查和算法公平。数据脱敏、访问控制、知情同意和合规审查能够降低数据使用风险，但并不能替代医学责任本身。医学人工智能只有在尊重患者权益、保障数据安全和接受临床复核的前提下，才能更稳妥地服务于疾病识别、风险预测和健康管理。







\subsection*{本节小结}

本节围绕医学数据的类型、处理与应用展开。电子病历、医学影像、生理信号和生命组学数据从不同层面记录人体健康状态，并共同表现出多源异构、专业性强、动态变化和高度敏感等特征。原始数据需要经过采集、清洗、标注和标准化处理，再通过数据对齐、特征表示与多源融合，才能被模型稳定利用。在此基础上，医疗人工智能可承担病灶识别、辅助诊断、风险预测和治疗决策支持等典型任务。同时，医学数据涉及患者隐私与诊疗安全，其采集和使用必须遵循数据脱敏、授权管理、伦理审查和算法公平等规范。理解数据类型、处理流程与医学任务之间的对应关系，是开展医学人工智能应用的基础。

\subsection*{本节习题}
\begin{enumerate}
\item 举例说明电子病历、医学影像、生理信号和生命组学数据在结构形式和医学用途上的主要差异。
\item 数据清洗与数据标准化的侧重点有何不同？请各举一例加以说明。
\item 结合具体病例，说明为什么单一来源的数据往往不足以支持疾病判断，多源数据融合能够提供哪些补充信息。
\item 病灶识别、辅助诊断、风险预测和治疗决策支持四类任务分别关注什么？它们对数据有哪些不同要求？
\item 使用医学数据训练模型时，应从哪些方面保护患者隐私并满足伦理规范？
\end{enumerate}

\clearpage
\section{医学影像识别与辅助诊断}

\begin{sectionkeypoints}
\item 了解X线、CT、MRI、超声、核医学影像和数字病理图像的成像特点，认识不同影像在结构显示、功能观察和病灶分析中的适用场景。
\item 掌握医学影像人工智能处理的一般流程，包括图像预处理、区域提取、图像分割、特征表示、模型训练和性能评价。
\item 理解肺结节检测、眼底病变筛查、皮肤病变识别、病理肿瘤分级和心脏影像分析等典型应用中的任务目标与方法选择。
\item 熟悉辅助诊断结果的主要表达方式，包括定位、分割、分类、定量和提示性结果，能够区分模型输出与临床诊断结论。
\item 认识医学影像辅助诊断的使用边界，理解图像质量、设备差异、病灶形态和患者临床资料对结果判读的影响。
\end{sectionkeypoints}
清晨的影像科阅片室里，检查列表已经排满屏幕，一次胸部CT检查可能包含数百张连续切片，医生需要逐层寻找隐藏在血管、骨骼和正常组织之间的细微异常。几毫米的肺结节、边界模糊的脑出血区域或眼底图像中的微小出血点，都可能影响后续诊疗判断。随着检查数量增加，人工阅片容易受到工作负荷、观察经验和病灶复杂程度的影响，医学影像人工智能由此被用于辅助发现异常、测量病灶和分析影像特征。

医学影像进入模型后，通常需要依次经过图像预处理、特征提取、异常识别、定量分析和结果展示等环节。图像预处理用于调整尺寸、灰度和清晰度，特征提取用于识别与组织结构及病灶形态相关的信息，模型随后可以通过边界框、分割轮廓、热力图或概率值提示可疑区域。例如，肺结节检测系统可以标出结节位置并测量直径，医生再结合结节密度、边缘形态和既往影像进行复核。如图\ref{fig:ch9_05}所示，医学影像人工智能并非跳过分析过程直接生成诊断结论，而是将图像逐步转化为可查看、可测量和可验证的辅助信息，为病灶识别与临床判断提供参考。


\par
\noindent\begin{minipage}{\textwidth}
\centering
\BookFigureImage{./images_chapter_07-08/9.5.png}
\captionof{figure}{医学影像识别与辅助诊断流程示意图}
\label{fig:ch9_05}
\end{minipage}
\par

\subsection{医学影像类型与成像特点}


同一处病变在不同医学影像中可能呈现出完全不同的样貌。骨折在X线影像中表现为骨皮质连续性中断，肺部结节在CT切片中显示为局部密度异常，脑组织病变在MRI中则可能呈现不同的信号变化；超声能够记录心脏瓣膜和胎儿等结构的实时运动，核医学影像显示组织代谢与功能活动，数字病理图像则把观察范围深入到细胞和组织结构层面。不同成像技术提供的信息各有侧重，需要结合检查目的和病变特点进行选择。

医学影像在空间分辨率、组织对比度、时间连续性和功能信息等方面存在差异，这些特点也会影响人工智能的处理方式。CT和MRI通常由连续切片组成，模型需要分析病灶在多个层面中的空间关系；超声图像容易受到探头角度、操作方式和组织运动影响，处理时需要关注图像质量与动态变化；数字病理切片范围大、细节多，模型常从局部图像块中识别细胞形态和组织结构。例如，在肿瘤评估中，CT可用于观察病灶大小与位置，核医学影像能够补充代谢活跃程度，病理图像则用于分析细胞和组织特征，三类信息分别对应宏观结构、功能状态和微观形态。

如图\ref{fig:ch9_06}所示，X线影像、CT、MRI、超声影像、核医学影像和数字病理图像分别从结构、动态、功能和微观层面提供医学信息。人工智能模型需要根据不同影像的成像特点完成质量控制、特征提取、病灶识别和定量分析，不能将适用于某一类影像的处理方法直接套用于其他影像类型。


\par
\noindent\begin{minipage}{\textwidth}
\centering
\BookFigureImage{./images_chapter_07-08/9.6.png}
\captionof{figure}{常见医学影像类型与成像特点示意图}
\label{fig:ch9_06}
\end{minipage}
\par

\subsubsection*{\textbf{•X线影像}}

急诊患者因摔伤出现小腿剧烈疼痛时，医生通常会先安排X线检查，骨皮质是否连续、骨端是否移位以及关节面是否受损，都能够在影像中较快显示。X线影像利用不同组织对X射线吸收程度的差异形成二维投影，高密度的骨骼吸收射线较多，在图像中通常呈现较亮区域；含气较多的肺组织吸收较少，通常呈现较暗区域。软组织的密度差异相对较小，彼此之间的对比通常不如骨骼和肺部明显。

临床常见的X线检查包括胸部X线片、骨关节X线片、腹部平片和乳腺X线摄影。胸部X线片可以观察肺野、心影和胸廓结构，常用于肺部感染、胸腔积液和部分心肺异常的初步筛查；骨关节X线片主要用于判断骨折、脱位和关节退行性改变；腹部平片可辅助观察肠梗阻、部分结石和消化道穿孔征象；乳腺X线摄影则通过显示肿块、结构扭曲和微小钙化等变化，辅助乳腺疾病筛查。

X线检查速度较快、操作方便且设备覆盖范围广，适合急诊检查、常规体检和基层医疗等场景。胸痛患者到达急诊后，胸部X线片可以较快提供肺部和胸腔的初步信息；骨折患者在复位或手术后，也可以通过多次X线检查观察骨端位置和愈合情况。检查结果还可帮助医生判断是否需要进一步进行CT或MRI检查，从而形成由初步筛查到深入评估的影像检查路径。

由于X线影像将三维人体结构投影到二维平面，不同组织可能相互重叠，较小病灶、低对比度异常和深部结构变化容易被遮挡。肺部结节可能与血管或肋骨影重合，早期骨折也可能因裂纹细小或投照角度不合适而难以发现。图像质量还会受到患者体位、呼吸配合、曝光条件和设备参数影响，因此判读时需要结合多角度影像、临床表现和其他检查结果。人工智能处理X线影像时，也需要关注组织重叠、图像对比度和设备差异，避免将正常解剖结构或成像伪影误判为病灶。

\subsubsection*{\textbf{•CT}}

一名突发剧烈头痛并伴有肢体无力的患者被送入急诊后，医生通常会迅速安排颅脑CT检查，通过连续断层图像判断是否存在脑出血及其范围。CT利用X射线从多个角度扫描人体，再由计算机重建横断面图像，减少了普通X线影像中不同组织相互重叠的影响。连续切片还可以进一步重建为冠状位、矢状位或三维图像，使器官、病灶及周围结构之间的空间关系更加清楚。

CT检查具有扫描速度快、空间分辨率较高和结构显示清晰等特点，常用于颅脑、胸部、腹部和骨关节等部位。颅脑CT可以快速识别出血、骨折和明显占位，胸部CT能够显示肺实质、气道、胸膜及纵隔结构，腹部CT可用于观察肝、胆、胰、脾和肾等器官。复杂骨折患者接受CT检查后，医生还可以通过三维重建观察骨折线、碎骨片和关节面的空间位置，为手术方案制定提供参考。

不同组织对X射线的吸收程度不同，因此空气、脂肪、软组织、骨骼和钙化在CT图像中呈现不同密度。肺部结节筛查中，薄层CT能够显示结节大小、密度和边缘形态；增强CT通过注射对比剂观察组织和血管的强化变化，可以辅助判断肿瘤血供、血管病变及病灶与周围组织的关系。图像层厚、重建方式和扫描参数会影响病灶显示效果，较厚切片可能掩盖微小结节，图像噪声过大也会干扰边界判断。

CT检查需要使用电离辐射，因此儿童、孕妇及需要多次复查的患者应根据病情控制检查频率和扫描范围。增强CT还需要关注对比剂过敏和肾功能情况，部分软组织病变的显示能力也不如MRI。人工智能分析CT图像时，需要结合连续切片之间的空间关系，并注意层厚、设备和重建参数差异，避免将血管截面、正常组织或成像伪影误认为病灶。


\subsubsection*{\textbf{•MRI}}

一名疑似脑卒中的患者进入MRI检查室后，同一处脑组织会在不同序列中呈现出不同信号，医生需要结合这些图像判断病灶位置和范围。MRI利用强磁场和射频脉冲激发人体组织中的氢质子，再根据产生的磁共振信号重建图像。由于不同组织的含水量、质子密度和弛豫特性存在差异，肌肉、脂肪、神经和实质器官能够呈现较清楚的对比。

MRI具有较强的软组织分辨能力，可直接获得横断面、冠状面和矢状面图像，常用于颅脑、脊柱、关节、腹部和肿瘤检查。膝关节损伤时，MRI可以显示半月板、韧带和软骨变化；颅脑检查中，则有助于观察梗死、肿瘤、炎症和脱髓鞘病变。增强MRI还能够显示病灶的血供、边界和强化方式，为病变范围判断提供补充信息。

MRI可以通过不同成像序列反映组织的多方面特征，其中，T1WI（T1-weighted imaging，T1加权成像）主要显示解剖结构，T2WI（T2-weighted imaging，T2加权成像）对液体和水肿较敏感，DWI（diffusion-weighted imaging，弥散加权成像）常用于观察水分子扩散受限，FLAIR（fluid-attenuated inversion recovery，液体衰减反转恢复序列）有助于显示脑实质异常。急性脑梗死在DWI中可能较早出现异常信号，因此临床判读通常需要结合多个序列，而不能只依赖单一图像。

MRI不使用电离辐射，但检查时间较长，患者移动、呼吸和金属植入物都会影响检查安全与图像质量，其对骨皮质和钙化的显示也通常不如CT直观。人工智能分析MRI时，常需完成图像配准、多序列融合、组织分割和病灶识别。脑肿瘤分析中，模型可以联合T1WI、T2WI、FLAIR和增强序列区分肿瘤核心、水肿区和周围组织，同时还需处理序列缺失、位置偏移和扫描参数差异等问题。

\subsubsection*{\textbf{•超声影像}}

心脏超声检查时，探头贴近患者胸壁，屏幕上的心腔、瓣膜和心肌随着心跳连续运动，医生可以同步观察结构变化和血流状态。超声成像利用探头发射高频声波，并接收声波在不同组织界面形成的回声，再经过信号处理生成二维灰度图像、动态视频或血流图像。二维超声主要显示器官形态和组织结构，彩色多普勒超声可以观察血流方向与速度，超声弹性成像则用于评估组织硬度和弹性差异。

超声检查具有实时、无创、便于重复检查和可床旁操作等特点，常用于心脏、腹部、甲状腺、乳腺、妇产科和血管检查。心脏超声能够观察心腔大小、瓣膜活动和心功能，腹部超声可显示肝、胆、胰、脾和肾等器官，血管超声则可辅助发现血管狭窄和血栓。急诊患者不便移动时，医生还可以在床旁快速观察心包积液、胸腔积液或腹腔内异常情况。

超声图像容易受到探头角度、操作者经验、患者体型和声窗条件影响，同一部位在不同切面下可能呈现不同形态。肠道气体和骨骼会阻挡声波传播，使部分深部结构难以显示；组织边界、回声强度和图像纹理还可能受到散斑噪声干扰。甲状腺结节检查中，探头压力和切面变化会影响结节形态与测量结果，因此判读时需要结合多个切面、动态表现和临床资料。

人工智能处理超声影像时，常用于标准切面识别、图像质量评价、器官与病灶分割、自动测量及动态功能分析。胎儿超声可以通过模型识别标准检查切面，心脏超声可以自动勾画心腔并计算相关功能参数，甲状腺和乳腺超声则可辅助分析结节边界、形态和回声特征。由于超声图像具有明显的操作者依赖性，模型训练和应用时还需要处理切面差异、图像噪声和设备变化带来的影响。

\subsubsection*{\textbf{•核医学影像}}

肿瘤患者完成PET检查后，图像中代谢活跃的区域会呈现较明显的示踪剂摄取，医生可以据此观察病灶分布及其功能状态。核医学影像通过向人体引入放射性核素示踪剂，记录其在组织和器官中的分布，再根据信号重建图像。与主要显示解剖结构的X线、CT和MRI相比，这类影像更侧重反映组织代谢、血流灌注和器官功能变化。

PET常用于肿瘤定位、分期和疗效评估，也可辅助分析神经系统及心血管疾病；SPECT则常用于心肌灌注、脑血流和骨显像等检查。肿瘤组织代谢活跃时，PET图像可能在结构变化尚不明显的阶段显示异常摄取；心肌灌注检查中，SPECT可以呈现不同区域的血流分布，为缺血范围判断提供参考。

核医学影像的空间分辨率通常低于CT和MRI，异常摄取区域的边界和解剖位置可能不够清楚，因此临床上常采用PET-CT、SPECT-CT或PET-MRI进行联合分析。CT或MRI提供器官结构和病灶位置，核医学影像补充代谢与功能信息，例如PET显示肺部存在异常摄取后，CT能够进一步确定该区域位于肺结节、淋巴结还是邻近组织。

人工智能处理核医学影像时，常用于图像配准、异常摄取区域识别、功能参数提取和联合影像分析。模型可以比较病灶治疗前后的摄取强度和范围变化，也可以结合CT或MRI完成病灶定位与风险判断。炎症、感染和正常生理活动也可能出现较高摄取，因此模型结果仍需结合解剖位置、患者病史和其他检查资料进行复核。


\subsubsection*{\textbf{•数字病理图像}}

病理医生打开一张肿瘤组织切片时，需要先在低倍视野中观察组织整体结构，再不断放大可疑区域，检查细胞核形态、排列方式和染色变化。数字病理图像通过对染色后的玻璃切片进行高分辨率扫描，将显微镜下的组织和细胞转化为可存储、放大和计算的图像数据。常见类型包括苏木精-伊红染色图像、免疫组织化学染色图像和其他特殊染色图像。

数字病理图像能够呈现细胞核大小与形态、细胞排列、组织结构和染色强度等微观特征，常用于肿瘤识别、病理分型、分级判断和组织学评估。一张全切片图像通常包含大量组织区域，低倍图像可用于观察肿瘤分布和组织结构，高倍图像则用于分析细胞异型性、核分裂象等细节。例如，乳腺肿瘤评估既需要判断病变区域范围，也需要观察肿瘤细胞形态和特定标志物的染色表达。

切片制作、染色批次和扫描设备会影响图像颜色、清晰度及背景噪声，同一种组织在不同实验室中可能呈现不同色调。切片折叠、组织缺失和扫描模糊也可能遮挡关键区域，炎症组织与肿瘤组织在局部形态上还可能较为接近。因此，病理判读需要综合组织结构、细胞形态、染色结果和临床资料，不能只依赖单个局部区域。

人工智能分析数字病理图像时，通常先将尺寸较大的全切片图像划分为多个局部图像块，再完成颜色标准化、组织区域识别、细胞核检测、病灶定位和分类分级。模型可以在局部图像中识别异常细胞，并将不同区域的分析结果汇总到整张切片，例如标出肿瘤分布范围、统计阳性细胞比例或辅助判断病理等级。由于局部切块可能割裂组织之间的空间关系，分析时还需要兼顾细胞细节与整体组织结构。
\subsubsection*{\textbf{•小结}}

本节介绍了X线、CT、MRI、超声、核医学和数字病理图像的成像特点、主要用途及人工智能处理方式。不同影像分别侧重组织结构、动态变化、代谢功能或细胞形态等信息，人工智能可用于图像质量评价、病灶识别、区域分割、自动测量和联合分析。模型结果容易受到成像设备、采集条件和图像质量等因素影响，仍需结合原始影像与临床资料进行判断。


\subsection{图像处理流程与识别方法}

一张刚从影像设备中导出的检查图像，可能同时存在噪声、运动伪影、灰度差异和扫描参数变化。胸部CT中的微小结节可能被血管影遮挡，超声图像中的组织边界也可能受到散斑噪声干扰。为了使模型获得稳定、清晰且具有医学意义的信息，原始影像通常需要经过图像预处理、区域提取与图像分割、特征表示、模型训练和结果评价五个环节。

图像预处理主要用于校正灰度差异、抑制噪声并减少伪影影响；区域提取与图像分割用于确定器官或病灶范围，例如从颅脑MRI中分离脑组织，或在胸部CT中勾画肺结节边界；特征表示进一步提取病灶的形态、纹理和空间结构，模型训练则利用已标注影像学习这些特征与疾病类别之间的关系。完成训练后，还需要通过独立数据评价病灶识别、分割和分类结果，检查模型是否存在漏检、误检或边界偏差。

如图\ref{fig:ch9_07}所示，各处理环节前后衔接，最终将原始影像转化为病灶位置、分割轮廓、类别概率或定量测量结果。肺结节识别系统可以标出可疑区域并测量结节大小，脑肿瘤分割模型可以显示病灶范围及其与周围组织的关系，医生则根据原始图像和临床资料复核模型输出，使人工智能结果成为可检查、可比较的辅助信息。


\par
\noindent\begin{minipage}{\textwidth}
\centering
\BookFigureImage{./images_chapter_07-08/9.7.png}
\captionof{figure}{医学影像处理流程与识别方法示意图}
\label{fig:ch9_07}
\end{minipage}
\par

\subsubsection*{\textbf{•图像预处理}}

两名患者的胸部CT中可能存在大小相近的肺结节，但由于检查设备、扫描层厚、重建算法和呼吸状态不同，结节在图像中的清晰度、边缘和灰度表现可能并不一致。模型若直接使用这些原始图像，可能把设备差异、扫描噪声或显示条件当作疾病特征，因此需要先对图像质量和空间表达进行调整。常见处理包括图像去噪、灰度归一化、对比度增强、尺寸调整与重采样、图像配准以及感兴趣区域提取。

图像去噪和增强主要用于减少采集干扰，并突出与当前任务相关的结构。低剂量CT能够降低辐射暴露，但图像中的颗粒噪声较明显，适度去噪有助于显示小结节和低对比病灶；超声图像中的散斑噪声可能掩盖组织边界，也需要在保留真实回声特征的基础上进行抑制。眼底图像经过对比度增强后，血管、出血点和渗出区域会更加清楚，但处理强度过大可能放大噪声或改变细小结构，因此不能单纯追求图像视觉上的清晰。

灰度与颜色归一化用于减小设备和采集条件造成的整体差异，使模型更关注病灶形态和组织结构。CT图像可以根据分析对象选择适当的窗宽窗位，使肺组织、骨骼或软组织得到更合适的显示；MRI在不同设备和序列下可能出现明显的信号强度差异，需要调整其数值分布。数字病理图像还常受到染色批次影响，同一种组织在不同切片中可能呈现深浅不同的颜色，通过染色归一化可以减少这种变化对细胞和组织识别的干扰。

尺寸调整与重采样主要用于统一图像的像素间距、层厚和体素大小，使不同病例中的空间尺度具有可比性。肺结节检测时，若一组CT采用1mm层厚，另一组采用5mm层厚，同样大小的结节在切片中的表现会明显不同；重采样可以将图像调整到相对统一的体素尺度，便于模型比较病灶大小和空间位置。调整过程中还需要避免过度插值，以免使微小病灶边缘变得模糊。

图像配准通过空间变换，使不同时间点、不同序列或不同成像方式中的同一解剖结构尽量重合。患者治疗前后的肿瘤影像经过配准后，可以更准确地比较病灶范围变化；颅脑MRI中的T1WI、T2WI和增强序列完成对齐后，模型能够联合分析各序列提供的信息；PET-CT分析中，代谢异常区域也需要与CT中的解剖位置准确对应。配准出现偏差时，病灶可能被错误映射到邻近组织，因此仍需检查关键结构是否对齐。

感兴趣区域提取可以减少无关背景对模型的干扰，使分析集中在目标器官或病变区域。肺结节检测常先分离肺野，排除胸壁和部分纵隔结构；肝肿瘤分析可以先定位肝脏区域，再识别内部病灶；眼底图像分析则可定位视盘和黄斑，为病变位置判断提供参考。若感兴趣区域划分过小，边缘病灶可能被遗漏，因此提取范围需要结合具体任务确定。

不同影像和任务需要采用不同的预处理组合，CT主要关注层厚、体素尺度和窗宽窗位，MRI更重视信号归一化和多序列配准，超声需要处理散斑噪声与切面差异，数字病理则侧重染色归一化、图像切块和组织区域筛选。所有处理步骤都应保留参数和原始图像，便于检查模型结果是否受到预处理影响，并避免在改善图像质量的同时改变具有医学意义的细节。
\subsubsection*{\textbf{•区域提取与图像分割}}

腹部CT中同时包含肝脏、脾脏、胃肠道、骨骼和脂肪等多种结构，肝肿瘤可能只占其中很小的区域。模型若直接在整幅图像中搜索病灶，容易受到邻近器官和背景组织干扰，因此通常先确定肝脏所在范围，再进一步勾画肿瘤边界。区域提取用于选出与当前任务相关的器官、组织或局部范围，图像分割则在此基础上将目标从背景中精确划分出来。

分割结果通常以二维图像中的像素或三维影像中的体素为单位表示，既能够显示病灶位置，也可以支持面积、体积和形态测量。脑出血患者复查CT时，医生可以比较两次血肿分割结果，判断出血范围是否扩大；肿瘤治疗前后，病灶体积变化可用于辅助评价治疗反应；心脏影像中的心腔分割还可用于计算心室容积和射血分数。边界出现偏差时，这些定量结果也会随之变化，因此分割精度会直接影响后续分析。

传统分割方法主要利用图像的灰度、边缘和区域相似性。阈值分割按照灰度或信号强度区分目标与背景，适合骨骼、肺野等对比较明显的结构；区域生长从指定位置出发，将特征相近的邻近区域逐步合并；边缘检测和活动轮廓模型则通过寻找灰度变化或调整轮廓位置贴近目标边界。脑出血区域与周围组织密度差异明显时，阈值方法可能取得较好效果，但面对边界模糊的浸润性肿瘤或灰度不均的MRI图像时，这些方法容易受到噪声和参数设置影响。

深度学习分割方法能够从标注图像中自动学习器官和病灶特征。FCN（fully convolutional network，全卷积网络）可以对图像中的每个位置进行类别判断，U-Net则通过编码器提取多层特征，再由解码器恢复空间细节，并利用跳跃连接保留目标边界。肺野、肝脏、脑肿瘤和心腔等任务中，U-Net及其改进结构应用较多，尤其适合医学影像标注样本相对有限的场景。

小病灶、多个目标同时出现或边界形态复杂时，还可以结合多尺度特征、注意力机制和实例分割方法。多尺度特征能够同时分析器官整体结构和局部细节，注意力机制可以减弱背景干扰，Mask R-CNN（mask region-based convolutional neural network，掩膜区域卷积神经网络）则能够分别定位并分割图像中的多个独立目标。数字病理图像中同时存在大量细胞核时，实例分割可以为每个细胞生成独立轮廓；眼底图像中的微小出血点和胸部CT中的小结节，也需要模型保留较精细的局部信息。

CT和MRI通常由连续切片组成，三维分割可以利用相邻层面之间的空间连续性。3D CNN（three-dimensional convolutional neural network，三维卷积神经网络）和3D U-Net能够直接处理体数据，使脑肿瘤、肺结节和肝脏的分割结果在多个切片之间保持连贯。二维方法计算量较小，但可能忽略前后切片关系；三维方法能够获得更完整的空间信息，却需要更多显存、训练数据和三维标注。

医学影像中的病灶边界并不总是清楚，炎症浸润区、早期肿瘤和低对比组织常与周围结构自然过渡，不同医生对边界范围也可能给出不同判断。扫描层厚、运动伪影和图像噪声还会进一步影响分割结果，因此手术规划、放疗靶区勾画和疗效评估等任务不能直接采用模型输出。医生需要返回原始图像检查分割边界，并结合解剖结构和临床资料进行必要修正。
\subsubsection*{\textbf{•特征提取与特征表示}}

医生观察胸部CT中的肺结节时，会综合判断其大小、密度、边缘毛刺和内部结构；查看皮肤镜图像时，则会关注颜色分布、边界规则性和病灶对称性。人工智能模型需要把这些视觉线索转化为数值、向量或特征图，才能完成病灶识别和疾病分类。特征提取用于从图像中获得与任务相关的信息，特征表示则将这些信息整理为模型能够计算和比较的形式。

传统方法通常根据医学知识预先设计灰度、纹理、形状和边缘等特征。灰度特征可以描述CT病灶的密度分布，纹理特征用于反映组织内部是否均匀，形状特征则可表示病灶大小、圆整程度和边界规则性。肺结节分析中常提取直径、体积、密度和毛刺等特征，病理图像分析则可计算细胞核大小、核质比、染色强度及组织排列特征，再将这些变量输入SVM、RF或LR等模型进行分类。此类特征医学含义较清楚，但对复杂背景和细微形态变化的表达能力有限。

影像组学可以从病灶区域中批量提取形状、灰度和纹理等定量特征，使影像信息由肉眼描述转化为可计算变量。肿瘤影像中可能一次得到数百甚至上千个特征，其中部分变量相互重复或与任务关系较弱，因此还需要进行特征筛选或降维。PCA可以将相关特征压缩为少量综合变量，特征选择方法则保留与良恶性判断或疗效预测关系较大的指标，从而减少小样本建模中的过拟合风险。

CNN能够从原始图像中自动学习不同层级的特征，浅层通常识别边缘、纹理和局部对比度，较深层则逐渐形成病灶形态和组织结构的综合表达。皮肤病变识别模型可以学习颜色不均和边界变化，眼底筛查模型能够关注血管走行、出血点和渗出区域，病理模型则可提取细胞核形态与腺体结构。模型输出的特征图保留一定的空间位置信息，适合定位和分割任务，压缩后的特征向量则常用于疾病分类、风险评分和疗效预测。

复杂影像往往同时包含整体结构与局部细节，因此还需要结合多尺度特征。肺结节检测既要识别几毫米的局部异常，也要分析结节与血管及肺野的空间关系；数字病理分析既要观察单个细胞核，也要判断腺体和组织的排列方式。通过融合不同分辨率或不同网络层级的特征，模型可以在局部细节和整体背景之间建立联系，减少微小病灶被忽略的情况。

注意力机制可以提高关键区域在特征表示中的权重，使模型减少对无关背景的关注。皮肤镜图像中的尺标、气泡和正常皮肤不应成为疾病判断依据，病理图像中的大片空白区域也不需要与肿瘤组织具有相同权重。热力图或注意力图还可以显示模型重点分析的位置，帮助医生检查其是否真正关注病灶，而不是依赖图像标记或设备特征作出判断。

医学影像标注数量有限时，还可以使用预训练模型和迁移学习，将模型在大规模数据中获得的基础图像特征应用于新的医学任务，再利用眼底、皮肤镜或胸部影像数据进行调整。由于自然图像与医学影像在结构和医学含义上存在差异，预训练模型仍需在目标数据中验证，不能直接视为适用于所有任务。

特征表示应尽量保留病灶形态、组织纹理和空间分布等稳定信息，同时减少设备型号、扫描参数和图像标记带来的干扰。若模型主要识别某家医院的图像风格，而不是疾病本身，它在更换设备或医院后便可能出现性能下降。因此，特征分析还需要结合外部数据和可视化结果，检查模型学习到的内容是否与医学任务一致。

\subsubsection*{\textbf{•识别方法与模型训练}}

肺结节筛查系统需要判断可疑阴影是否属于结节，眼底筛查系统需要识别出血、渗出等病变，皮肤镜图像分析还要区分色素痣、基底细胞癌和黑色素瘤等不同类型。医学影像识别就是根据图像特征判断目标、类别或病变状态，其中，整幅图像的正常与异常判断属于分类任务，病灶位置提示属于目标检测任务，病灶边界勾画则需要结合图像分割方法。

传统机器学习方法通常先提取灰度、纹理、形状和边缘等人工特征，再使用LR、SVM、RF、KNN（k-nearest neighbors，K近邻）或GBDT完成分类。肺结节分析可以输入直径、密度、边缘毛刺和纹理等变量，乳腺影像分析则可利用肿块形态与钙化分布判断良恶性。此类方法适合样本量较小、特征含义明确的任务，医生也较容易理解模型依据，但其识别能力会受到人工特征设计范围的限制。

深度学习模型可以直接从原始影像中自动学习特征，CNN及其不同结构已广泛用于医学图像分类和病灶识别。VGG（Visual Geometry Group network，视觉几何组网络）采用较规整的卷积结构，ResNet（residual network，残差网络）通过残差连接支持更深层的特征学习，DenseNet（densely connected convolutional network，密集连接卷积网络）加强不同网络层之间的信息传递，EfficientNet则兼顾识别性能和计算量。胸部X线分类中，模型可以逐步学习肺纹理和局部阴影；皮肤病变识别中，则能够综合颜色、边界和局部结构形成类别判断。

不同影像的数据结构决定了模型的处理方式。胸部X线、眼底照片和皮肤镜图像通常采用二维模型进行分析，CT和MRI则由连续切片组成，单张切片可能无法完整呈现病灶形态，因此可使用3D CNN或3D U-Net处理三维空间信息。脑肿瘤分析中，三维模型能够结合相邻切片观察病灶范围；多序列MRI或PET-CT任务还可以设置不同分支分别提取特征，再在模型内部完成联合分析。

模型训练前需要将数据划分为训练集、验证集和测试集，三者分别用于学习模型参数、调整训练设置和评价模型在未见数据上的表现。划分时应避免同一名患者的不同切片同时出现在训练集和测试集中，否则模型可能通过识别患者或检查特征获得虚高结果。标签也需要与任务保持一致，肺结节模型可以使用“良性/恶性”标签，眼底筛查模型可以采用不同病变等级，分割模型则需要病灶或器官的边界标注。

训练过程中，损失函数用于衡量预测结果与真实标签之间的差异，优化方法根据这一差异调整模型参数。分类任务常使用交叉熵损失，分割任务可以结合Dice损失，使模型生成的区域更接近医生标注。模型在训练图像上反复预测并修正参数后，还需要观察验证集表现；若训练损失持续下降，而验证性能开始变差，通常说明模型出现了过拟合。

医学影像数据常面临样本不足和类别不平衡问题，可以通过数据增强、类别权重调整、CV（cross-validation，交叉验证）、正则化和早停等方法改善训练。图像旋转、翻转、裁剪和亮度调整能够增加样本变化，但增强方式应符合医学结构，例如带有明确左右方向的影像不能随意翻转。罕见病灶样本数量较少时，还需要提高其训练权重或补充病例，避免模型因正常图像数量过多而倾向于输出阴性结果。
分类模型通常采用准确率、敏感度、特异度、精确率、F1（F1 score，F1分数）和AUC（area under the receiver operating characteristic curve，受试者工作特征曲线下面积）进行评价，不同指标分别反映总体判断、阳性检出、阴性排除及分类平衡情况。
图像分割任务常使用DSC（Dice similarity coefficient，Dice相似系数）和IoU（intersection over union，交并比）评价预测区域与人工标注区域的重合程度。DSC表示预测区域与真实区域重叠部分的两倍，占两者面积或体积总和的比例；IoU表示两者交集占并集的比例。两项指标通常取0至1之间的数值，越接近1，说明模型勾画的区域与人工标注越一致。脑出血分割中，较高的DSC和IoU有助于提高血肿体积测量的准确性；肝肿瘤分割中，指标较低可能说明模型遗漏了部分病灶，或者把周围正常组织错误划入肿瘤区域；心腔分割若存在明显偏差，还会影响心室容积和射血分数的计算结果。无论采用哪类指标，都不能只依据单个数值判断模型性能，一个总体准确率较高的模型若频繁漏掉数量较少的高风险病灶，仍不适合直接用于临床筛查。

识别结果可以表现为类别概率、病灶位置框、分割轮廓或热力图，并应与原始影像建立明确对应。眼底模型给出病变等级时，可以同时标出出血和渗出区域；肺结节模型输出恶性概率时，也可以展示结节位置和重点关注区域。模型完成内部测试后，还需要在不同医院、设备和患者人群中进行外部验证，并由医生结合原始影像和临床资料复核结果，避免把模型输出直接作为最终诊断。


\subsubsection*{\textbf{•结果输出与性能评价}}

医生使用医学影像辅助诊断系统时，接触到的通常不是网络参数和计算过程，而是能够与原始影像对应的分析结果。肺结节检测系统可以在CT图像上标出可疑结节，并给出直径、体积和风险评分；眼底筛查系统可以提示糖尿病视网膜病变的分级，同时标记出血点和渗出区域；心脏影像分析系统则可显示心腔轮廓，并计算心室容积和射血分数。结果只有能够被医生理解、复核和追溯，才能用于后续诊疗分析。

模型输出形式应与任务目标保持一致，分类模型通常给出类别标签和概率值，如判断乳腺肿块为良性或恶性；检测模型通过定位框提示病灶的大致位置，如在胸部CT中标出多个肺结节；分割模型以轮廓或掩膜显示器官和病灶边界，如勾画脑肿瘤、水肿区或放疗靶区；定量分析模型则输出直径、面积、体积、密度和功能参数。热力图还可以显示模型重点关注的位置，帮助医生检查其判断是否基于病灶，而不是图像标记或无关背景。

不同临床任务需要选择不同的评价重点，疾病筛查更关注敏感度，因为漏掉真正患病者可能延误进一步检查；确诊辅助需要兼顾敏感度和特异度，避免大量误报增加患者负担；在类别分布不均衡的任务中，还应结合精确率、F1和AUC进行判断。罕见肺部病变识别中，即使模型总体准确率较高，只要大部分阳性病例仍被误判为正常，就不能说明其具有可靠的筛查能力。

图像分割任务主要评价预测区域与人工标注区域的一致程度，常使用Dice系数和交并比。脑出血分割中，较高的区域重合度有助于获得更准确的血肿体积；放疗靶区勾画还需要关注边界距离，因为局部轮廓偏移可能影响照射范围。目标检测任务除统计病灶检出率外，还应分析误报数量及不同大小病灶的识别效果，例如肺结节模型可能容易发现较大的实性结节，却漏掉微小磨玻璃结节。

模型在独立测试集中的表现，可以用于检验其是否真正学习了疾病相关特征，而不是仅仅记住训练样本中的图像模式；进一步采用外部验证，还可以观察模型在不同医院、设备和患者人群中的稳定性。同一肺结节模型从薄层CT数据应用到层厚较大的数据后，结节边缘和体积表现可能发生变化，眼底模型面对不同相机拍摄的图像时，也可能受到亮度、视野范围和成像质量差异影响。因此，仅在单一医院内部获得良好结果，还不足以说明模型具有广泛适用性。

图像亮度、裁剪范围或噪声发生轻微变化时，模型输出若出现明显波动，说明其对输入条件较为敏感；患者在短期内接受重复检查时，相近影像也应得到相对一致的判断。热力图、注意力区域和重要特征提示可以帮助医生检查模型依据，例如皮肤病变模型应主要关注病灶边界、颜色变化和内部结构，而不应把尺标、气泡或正常皮肤区域作为主要判断依据。结果稳定性与可解释性共同反映模型是否能够在临床环境中保持可靠表现。

医学影像模型进入临床流程后，仍需持续记录漏检、误报、医生修正和设备变化等情况。设备软件升级、扫描参数调整或患者来源改变，都可能导致原有模型性能下降。医院因此需要定期抽查模型结果，并保留原始影像、输出结果和修正记录，使模型在实际应用中始终处于可检查、可追溯和可控制的状态。

\subsubsection*{\textbf{•小结}}

本节介绍了医学影像识别的基本流程，包括图像预处理、目标区域提取、图像分割、特征提取、模型训练、结果输出和性能评价。图像预处理用于改善图像质量与数据一致性，目标定位和分割用于确定病灶范围，特征提取与模型训练用于完成异常识别或类别判断，性能评价则用于检验模型结果的准确性、稳定性和可靠性。不同影像类型和任务目标需要采用相应的处理流程与识别方法。

\subsection{典型疾病的影像识别应用}

医学影像中的异常表现会随疾病类型、成像方式和病变阶段发生变化，因此模型需要处理的目标也并不相同。胸部CT中的早期肺结节可能只有数毫米，还可能与血管截面或纤维条索相似，模型需要结合连续切片完成病灶检测、位置定位和风险判断；眼底图像中的微血管瘤、出血点和渗出区域大小不一、分布分散，分析时既要识别局部病变，也要评价视网膜病变程度。皮肤镜图像中的良性痣与黑色素瘤可能具有相近的颜色和形态，模型则需综合病灶对称性、边界规则性、颜色分布和内部结构进行分类。

部分疾病的影像分析还需要从病灶识别延伸到组织结构和功能状态评价。数字病理图像中的肿瘤分级需要同时观察细胞核异型性、腺体结构和组织排列，模型通常先定位可疑组织，再汇总不同局部区域的特征；心脏影像分析需要勾画心房、心室和心肌边界，并根据连续图像计算心室容积、室壁运动和射血分数。医学影像人工智能因而需要根据疾病的影像表现和临床目的，选择目标检测、图像分割、疾病分类或定量分析方法，并使模型输出能够回答病灶是否存在、位于何处、范围多大以及严重程度如何等具体问题。


\subsubsection*{\textbf{•肺结节检测}}

肺结节是肺部影像中出现的局灶性类圆形阴影，通常直径不超过3cm，其形成可能与炎症、良性增生、早期肺癌或其他肺部病变有关。不同肺结节在密度和形态上存在差异，实性结节内部密度较高，磨玻璃结节表现为轻度密度增高且仍可见内部血管，部分结节还可能出现毛刺、分叶、空泡或胸膜牵拉等征象。临床上，医生需要结合结节大小、密度、边缘形态、所在位置及随访变化判断其风险，而不能仅凭一次检查直接确定病变性质。

肺结节筛查主要采用LDCT（low-dose computed tomography，低剂量计算机体层成像）或常规胸部CT。一次检查通常包含数百张连续切片，直径仅为数毫米的结节可能隐藏在血管、支气管和正常肺纹理之间。血管截面在单张切片中也可能呈现圆形，局部炎症影和胸膜增厚有时与结节形态相近，微小磨玻璃结节还容易受到图像噪声干扰，因此医生需要连续观察前后切片，判断可疑阴影是否具有相对独立的立体结构。

人工智能检测通常先完成肺野分割，将肺组织从胸壁、纵隔和骨骼等结构中分离出来，再在肺野内部寻找候选区域。候选区域是模型初步筛出的可疑阴影，其中既可能包含真正的结节，也可能包含血管交叉、支气管壁或成像伪影。3D CNN可以同时分析相邻切片中的空间信息，提取结节的形态、密度及其与血管、支气管和胸膜之间的关系；多尺度特征融合则有助于同时识别较大的实性结节和边界较淡的微小磨玻璃结节。

如图\ref{fig:ch9_08}所示，肺结节人工智能检测依次经过影像输入、肺野分割、候选区域提取、3D CNN分析、结果输出和医生复核。系统可以标出结节位置，并给出直径、体积、密度类型和风险评分；患者复查时，还可比较结节是否增大、密度是否升高以及实性成分是否增加。模型结果只能作为筛查和辅助判断依据，医生仍需结合毛刺、分叶和胸膜牵拉等影像表现，以及患者年龄、吸烟史、随访结果和病理检查作出综合判断。

\par
\noindent\begin{minipage}{\textwidth}
\centering
\BookFigureImage{./images_chapter_07-08/9.8.png}
\captionof{figure}{肺结节人工智能检测流程示意图}
\label{fig:ch9_08}
\end{minipage}
\par

\subsubsection*{\textbf{•眼底病变筛查}}

眼底病变是发生在视网膜、视盘、黄斑和眼底血管等部位的异常改变，可能由糖尿病、高血压、年龄增长或其他眼科疾病引起。病变早期未必出现明显视力下降，但眼底已经可能出现血管扩张、出血、渗出或局部水肿。DR（diabetic retinopathy，糖尿病视网膜病变）是糖尿病常见的眼部并发症，病情进展后可能损伤黄斑或形成新生血管，严重时可导致视力下降甚至失明，因此早期筛查和定期随访具有重要意义。

眼底病变筛查主要使用眼底照相获得视网膜图像，图中可以观察视盘、黄斑和血管走行等结构。DR早期常出现微血管瘤，在图像中表现为细小红点；病情进一步发展后，还可能出现点状或片状出血、硬性渗出、棉絮斑和新生血管。高血压相关眼底改变则可能表现为血管变细、动静脉交叉异常或视盘水肿，年龄相关性黄斑病变主要累及黄斑区域，可能出现玻璃膜疣、色素紊乱或局部出血。不同病变在颜色、形态、大小和分布位置上存在差异，微小病灶还容易受到图像模糊、曝光不均和血管背景干扰。

人工智能筛查通常先对眼底图像进行亮度校正、视野裁剪和质量评价，再利用CNN提取血管、出血、渗出和黄斑区域等特征。部分系统采用迁移学习，在已有模型参数基础上继续使用眼底图像训练，以减少标注样本不足带来的影响；病灶检测模型还可以标出微血管瘤、出血点和渗出区域，使医生查看分级结果所依据的图像位置。图像中央存在硬性渗出时，模型需要判断病变是否靠近黄斑，因为相同面积的渗出出现在黄斑附近，往往具有更高的视功能风险。

如图\ref{fig:ch9_09}所示，眼底病变人工智能筛查通常经过图像输入与质量控制、病灶定位、特征提取、分类分级和转诊提示等步骤。系统可以输出是否存在病变、DR等级、可疑区域及是否建议转诊，帮助基层筛查人员优先识别需要进一步检查的患者。图像受到白内障、眼底反光或瞳孔过小影响时，模型可能无法清楚识别微小病灶；视盘附近的正常血管结构也可能被误判为出血，因此医生仍需结合患者病史、视力检查和眼科专科检查复核模型结果。

\par
\noindent\begin{minipage}{\textwidth}
\centering
\BookFigureImage{./images_chapter_07-08/9.9.png}
\captionof{figure}{眼底病变人工智能筛查流程示意图}
\label{fig:ch9_09}
\end{minipage}
\par


\subsubsection*{\textbf{•皮肤病变识别}}

皮肤病变是皮肤表面在颜色、形态、结构或生长方式上出现的异常改变，既包括色素痣和SK（seborrheic keratosis，脂溢性角化）等常见良性病变，也包括BCC（basal cell carcinoma，基底细胞癌）和黑色素瘤等恶性肿瘤。多数色素痣长期保持稳定，但黑色素瘤可能发生快速增大、颜色加深、边界改变或表面破溃，并具有浸润和转移风险。医生通常结合病变出现时间、近期变化、患者年龄、日晒史和家族史进行初步判断，必要时通过病理活检明确病变性质。

皮肤镜检查可以放大观察肉眼难以分辨的色素分布、边界形态、纹理结构和血管特征。良性色素痣常表现为形态较对称、边界相对规则和颜色分布较均匀，黑色素瘤则可能出现明显不对称、边界参差、多种颜色混合及异常色素网络。脂溢性角化图像中可能观察到类似角栓或囊状结构，BCC则常表现为树枝状血管、蓝灰色结构或局部溃疡。不同病变之间仍可能存在相似表现，早期黑色素瘤有时与普通色素痣接近，炎症和抓挠也可能改变病灶颜色与边界，因此仅凭单项特征难以完成可靠判断。

人工智能识别通常先分割病灶区域，减少正常皮肤和图像背景的干扰，再利用CNN或EfficientNet提取颜色、纹理、边界和结构等特征。病灶分割能够明确分析范围，避免模型把周围皮肤、尺标或检查标记当作诊断依据；多尺度特征可以同时观察整个病灶的对称性和局部细小结构，注意力机制则有助于提高不规则边缘、异常色素网络和血管区域在分类过程中的权重。黑色素瘤识别中，模型需要综合深浅不一的棕色、黑色或蓝灰色区域，而不能只因病灶颜色较深就直接判断为恶性。

如图\ref{fig:ch9_10}所示，皮肤病变识别通常经过皮肤镜图像输入、病灶分割、深度特征提取、注意力分析、类别判断和风险评估等步骤。系统可以输出病变类别、恶性风险概率和重点关注区域，为医生确定是否需要进一步检查提供参考。毛发遮挡、反光、气泡和图像边缘可能掩盖病变结构，颜色较浅的早期黑色素瘤也可能被误判为良性，因此模型结果不能替代皮肤科检查和病理诊断。医生仍需结合病变近期是否增大、是否出现瘙痒或出血，以及患者个人史和病理活检结果作出最终判断。

\par
\noindent\begin{minipage}{\textwidth}
\centering
\BookFigureImage{./images_chapter_07-08/9.10.png}
\captionof{figure}{皮肤病变识别流程与风险评估示意图}
\label{fig:ch9_10}
\end{minipage}
\par

\subsubsection*{\textbf{•病理肿瘤分级}}

病理肿瘤分级是根据肿瘤组织和细胞与正常组织的差异程度，评价肿瘤分化水平和生物学侵袭性的过程。分化程度较高的肿瘤通常仍保留部分正常组织结构，分化程度较低的肿瘤则可能出现明显的细胞异型、组织排列紊乱和核分裂增多，其生长和侵袭风险往往更高。病理医生需要综合观察细胞核大小与形态、腺体结构、组织排列、染色特点和核分裂象等信息，例如乳腺癌分级会关注腺管形成、细胞核异型性和核分裂活跃程度，前列腺癌则根据腺体结构模式进行Gleason评分，以辅助判断肿瘤恶性程度和治疗方案。

病理肿瘤分级主要使用数字化后的组织切片，尤其是WSI（whole slide image，全切片图像）。一张WSI可以同时包含正常组织、炎症区域、肿瘤组织、坏死区域和切片背景，图像尺寸远大于普通医学影像，模型通常无法直接一次处理整张切片。分析时需要先去除空白背景和明显伪影，再将组织区域划分为多个局部图像块。低倍图像可用于观察肿瘤分布和整体组织结构，高倍图像则用于分析细胞核、腺体边界和核分裂象，例如前列腺癌组织中的腺体融合、筛状结构和腺腔消失等变化，需要在不同放大倍数下综合判断。

人工智能分析通常结合图像切块、局部特征提取、多实例学习和注意力机制。MIL（multiple instance learning，多实例学习）是将一张切片中的多个局部图像块作为一个样本集合进行训练，使模型在只有切片级诊断或分级标签的情况下，学习哪些局部区域与肿瘤等级关系较大。注意力机制可以进一步为不同图像块分配权重，使含有典型肿瘤结构的区域对最终判断产生更大影响，而空白区域、正常脂肪组织和轻微炎症区域的影响相对减小。对于需要精细分析细胞形态的任务，还可以使用细胞核分割和目标检测方法，统计细胞核大小、形状、密度及核分裂象数量，并分析肿瘤细胞与周围间质、淋巴细胞和血管之间的空间关系。

如图\ref{fig:ch9_11}所示，数字病理肿瘤分级通常经过WSI输入、组织区域筛选、图像切块、局部特征提取、MIL分析、关键区域定位和分级结果输出等步骤。系统可以标出对分级贡献较大的组织区域，并提供肿瘤类型、等级或评分参考，帮助病理医生优先检查可疑区域。切片折叠、染色深浅差异、组织缺失和扫描模糊可能影响模型判断，肿瘤内部还可能同时存在不同分化程度的区域，因此模型输出不能替代病理医生的整体观察。最终分级仍需结合完整切片、免疫组织化学结果、临床资料及必要的分子检测进行确认。

\par
\noindent\begin{minipage}{\textwidth}
\centering
\BookFigureImage{./images_chapter_07-08/9.11.png}
\captionof{figure}{数字病理肿瘤分级流程示意图}
\label{fig:ch9_11}
\end{minipage}
\par


\subsubsection*{\textbf{•心脏影像分析}}

心脏疾病可引起心腔结构、心肌运动、瓣膜功能和泵血能力等方面的异常。心力衰竭是心脏无法满足机体血液循环需求的临床综合征，常表现为心室扩大、心肌收缩减弱或充盈功能异常；心肌病可导致心肌增厚、扩张或局部纤维化；瓣膜病会影响瓣膜开放与关闭，使血流出现狭窄或反流；心肌缺血则可能造成局部室壁运动减弱。医生需要通过影像检查观察这些结构和功能变化，并结合患者胸闷、气促、心悸等症状判断病情。

心脏影像主要包括心脏超声、心脏磁共振成像和心脏CT。心脏超声可以实时显示心腔、瓣膜和室壁运动，适合观察瓣膜反流、心室收缩和心包积液；心脏MRI能够显示心肌结构、运动及组织特征，可用于分析心肌水肿、瘢痕和纤维化；心脏CT则可清楚显示冠状动脉、心脏解剖结构和钙化情况。与主要关注静态病灶的肺结节检测不同，心脏影像分析还需要比较收缩期和舒张期的连续变化，判断心脏泵血过程是否正常。

人工智能分析通常先分割左心室、右心室、心肌和其他心腔结构，再根据连续帧计算面积、容积和运动变化。U-Net及其改进模型可以用于勾画心腔和心肌边界，时序模型则可分析心脏超声视频或动态MRI中的连续运动。系统需要识别舒张末期和收缩末期图像，其中舒张末期心室容积最大，收缩末期心室容积最小。根据两者可以计算EF（ejection fraction，射血分数），即心室每次收缩排出的血量占舒张末期心室总血量的比例，EF较低通常提示心脏收缩功能减弱。室壁运动分析还可发现某一区域收缩幅度不足，为心肌缺血或心肌梗死评估提供参考。

如图\ref{fig:ch9_12}所示，心脏影像分析通常经过影像输入、图像预处理、心腔与心肌分割、动态时序分析、定量参数计算和综合结果输出等步骤。系统可以显示心腔轮廓、心室容积、EF值和室壁运动异常区域，并比较患者治疗前后的心功能变化。心律失常、切面不标准、呼吸干扰和心内膜边界模糊都可能影响自动测量结果，例如心房颤动患者不同心动周期的心室容积可能存在明显差异，因此医生仍需结合原始影像、心电图、症状和其他检查结果复核模型输出。

\par
\noindent\begin{minipage}{\textwidth}
\centering
\BookFigureImage{./images_chapter_07-08/9.12.png}
\captionof{figure}{心脏影像分析流程与心功能评估示意图}
\label{fig:ch9_12}
\end{minipage}
\par

\subsubsection*{\textbf{•小结}}

本节介绍了肺结节检测、眼底病变筛查、皮肤病变识别、病理肿瘤分级和心脏影像分析等典型应用。针对不同疾病的影像特点，人工智能可采用3D CNN、CNN、EfficientNet、MIL、U-Net及注意力机制等方法，完成病灶定位、图像分类、区域分割和定量分析。模型输出可为疾病筛查与诊断提供参考，但其可靠性会受到图像质量、样本差异和临床信息完整性的影响，实际应用中仍需由医生结合患者资料进行复核。


\subsection{医学影像辅助诊断系统}

放射科医生面对一组包含数百张切片的胸部CT时，需要逐层寻找微小结节，并记录其大小、密度和位置；急诊医生查看颅脑影像时，则需要尽快发现出血、梗死或占位等异常。随着医学影像实现数字化存储和传输，图像中的灰度、形态、纹理和空间关系可以被计算机自动测量和分析，由此形成医学影像CAD（computer-aided diagnosis，计算机辅助诊断）系统。该系统利用图像处理、机器学习和深度学习等方法，对医学影像中的器官、组织和病灶进行检测、分割、分类及定量分析，并将分析结果呈现给医生作为诊断参考。

CAD的功能需要根据临床任务进行设置。胸部CT分析系统可以标出疑似肺结节，测量其直径和体积，并比较不同随访时间的变化；颅脑CT系统可以识别疑似出血区域，计算血肿体积并提示占位效应；心脏影像系统能够勾画心腔和心肌边界，计算射血分数等功能参数；眼底筛查系统则可识别微血管瘤、出血和渗出，并给出病变等级和转诊建议。医生通过这些定位框、分割轮廓、概率评分和定量结果，可以更快地找到可疑区域，也能够在复查过程中获得较一致的测量结果。

医学影像辅助诊断系统通常由影像数据接入、图像质量控制、预处理、器官与病灶识别、定量计算、结果展示和医生复核等部分组成。系统从DICOM格式的影像中读取图像及检查信息，经过灰度调整、重采样或图像配准后，再调用检测、分割或分类模型完成分析。肺结节系统需要结合连续CT切片判断可疑阴影是否属于立体结节，乳腺影像系统需要综合肿块形态和钙化分布，PET-CT系统则需要把代谢异常区域与具体解剖位置对应起来，因此不同系统虽然具有相似流程，内部模型和结果形式仍需与影像类型及疾病特点相匹配。

如图\ref{fig:ch9_13}所示，CAD将影像输入、图像处理、模型识别、定量分析和结果复核连接为连续流程，并可应用于肺结节检测、眼底病变筛查、心功能评价和肿瘤影像分析等任务。系统输出仍可能受到图像质量、扫描参数、病灶复杂程度和训练数据来源影响，例如血管截面可能被肺结节模型误报，运动伪影可能干扰心腔分割，眼底反光也可能遮挡微小病灶。因此，医生仍需结合原始影像、患者症状、实验室检查和病理结果进行判断，CAD承担的是信息提示和定量辅助作用，而不是替代医生作出最终诊断。

\par
\noindent\begin{minipage}{\textwidth}
\centering
\BookFigureImage{./images_chapter_07-08/9.13.png}
\captionof{figure}{医学影像辅助诊断系统流程与疾病应用示意图}
\label{fig:ch9_13}
\end{minipage}
\par

\subsubsection*{\textbf{•影像数字化}}

医生在传统阅片灯前观察X线胶片时，只能依靠胶片上固定的亮度和对比度判断病变，难以自由放大局部区域，也不便同时比较患者多次检查结果。医学影像数字化是将人体组织产生的成像信号转化为数字图像，并通过计算机完成显示、存储、传输和处理的过程。随着CT、MRI、CR（computed radiography，计算机X射线摄影）和DR（digital radiography，数字X射线摄影）等技术的应用，医学影像逐渐由胶片记录转向数字化管理。

CR通过成像板接收X射线信息，再利用专门设备读取并转换为数字图像；DR则使用数字探测器直接获得影像，成像和传输速度通常更快。数字化后的胸部X线图像可以调节亮度和对比度，使肺纹理、骨骼或软组织得到更合适的显示；连续CT切片可以进行冠状面、矢状面和三维重建，帮助医生观察肿瘤与血管之间的空间关系；患者复查时，系统还可以并列显示新旧影像，便于比较肺结节大小或脑出血范围的变化。

数字化影像以像素值、体素值和空间坐标等形式保存组织信息，使计算机能够读取图像中的灰度、密度、信号强度和解剖位置。肺结节模型可以在连续CT切片中计算结节直径和体积，心脏MRI系统可以根据不同心动周期测量心室容积，眼底分析系统则能够定位微血管瘤和出血区域。影像数字化由此为病灶检测、图像分割、定量测量和人工智能识别提供了数据基础。

\subsubsection*{\textbf{•辅助诊断系统}}

当医生需要在大量影像中寻找微小异常时，计算机可以先完成初步筛查，再把可疑区域提示在阅片界面上。CADe（computer-aided detection，计算机辅助检测）主要回答病灶是否存在以及位于何处，常用于发现肺结节、乳腺钙化、骨折线和眼底出血点；CADx（computer-aided diagnosis，计算机辅助诊断）则在病灶已经发现的基础上，进一步分析其类别、等级或风险，例如判断乳腺肿块的良恶性倾向，或对糖尿病视网膜病变进行分级。两类系统可以单独使用，也可以组成从病灶发现到风险判断的连续流程。

医学影像辅助诊断系统通常包括影像接入、图像预处理、模型分析、结果展示和医生复核等环节。系统可从PACS（picture archiving and communication system，影像归档与通信系统）中读取患者影像，经过质量检查、灰度调整、重采样或图像配准后，再调用检测、分割和分类模型完成分析。胸部CT系统可以用定位框标出可疑结节，并给出大小和密度信息；急诊颅脑CT系统可以提示疑似出血区域并估算血肿体积；乳腺影像系统则可以标记成簇钙化或可疑肿块，帮助医生优先检查高风险区域。

模型分析结果通常以病灶定位框、分割轮廓、类别概率、风险等级、热力图或结构化报告等形式显示。医生可以返回原始图像检查系统提示是否符合解剖结构和病变表现，并决定接受、修改或忽略模型结果。辅助诊断系统承担的是信息筛选、定量测量和风险提示功能，最终诊断仍需结合患者症状、病史、实验室检查和其他影像资料完成。

\subsubsection*{\textbf{•核心技术}}

医学影像辅助诊断系统需要根据临床任务选择图像分割、目标检测、图像分类和定量测量等技术。图像分割用于勾画器官或病灶边界，肺野分割可以排除胸壁和纵隔对结节检测的干扰，心室分割则为心室容积和射血分数计算提供基础。目标检测用于寻找异常区域的位置，胸部CT中的肺结节、X线图像中的骨折线以及眼底图像中的出血点，都可以通过检测模型进行定位。图像分类进一步判断图像或局部区域所对应的疾病类别，例如区分良性色素痣与黑色素瘤，或判断病理组织的肿瘤等级。

定量测量能够将影像中的病变转化为直径、面积、体积、密度和功能参数等数值。患者接受肿瘤治疗后，系统可以比较治疗前后病灶体积是否缩小；脑出血患者复查时，可以判断血肿范围是否扩大；心脏影像中，模型还可以根据舒张末期和收缩末期心室容积计算射血分数。与单纯使用“增大”“减小”或“功能下降”等描述相比，定量结果更便于随访比较和疗效评价。

CNN可以从影像中自动学习边缘、纹理、形态和空间结构等特征，U-Net常用于器官与病灶分割，3D CNN则能够利用CT和MRI连续切片之间的空间联系。肺结节模型需要结合前后切片区分结节与血管截面，脑肿瘤模型需要利用三维信息保持病灶边界连续，动态心脏超声和内镜视频还可以结合时序分析方法，识别室壁运动、瓣膜活动或病灶随时间发生的变化。注意力机制可以提高关键区域在分析过程中的权重，但模型关注位置仍需通过热力图或结果叠加图进行检查。

\subsubsection*{\textbf{•临床应用流程}}

医学影像辅助诊断系统进入临床后，需要与检查、阅片、报告和复核等实际环节衔接，而不能作为独立于医院工作流程之外的分析工具。胸部CT完成扫描后，影像首先传入PACS，肺结节系统读取图像并检查切片厚度、扫描范围和图像质量，再完成肺野分割和候选结节检测。系统随后排除部分血管截面、支气管壁和成像伪影，在阅片界面中标出可疑结节，并给出直径、体积、密度类型和风险提示。医生查看模型标记后，还需结合结节是否存在毛刺、分叶或胸膜牵拉，以及患者吸烟史和既往影像作出判断。

心脏MRI分析通常需要先识别不同心动周期中的左心室、右心室和心肌边界，再确定舒张末期与收缩末期图像，计算心室容积和射血分数。扩张型心肌病患者可能出现心室扩大和收缩功能下降，肥厚型心肌病患者则可能出现心肌局部明显增厚，系统可以通过自动分割和测量减少重复勾画工作。心律失常、呼吸运动和图像切面偏差仍可能影响结果，因此医生需要检查自动生成的心腔轮廓，并对明显错误区域进行修正。

急诊颅脑CT系统也需要满足快速响应和结果优先显示的要求。当系统发现疑似脑出血时，可以在影像列表中提升该检查的处理优先级，并提示出血位置和范围，使医生尽快复核。真正用于临床的软件除需要获得稳定的内部测试结果外，还应完成独立验证、外部验证、质量控制和相应的合规审查，以确认系统能够在预定设备、患者人群和诊疗场景中可靠运行。

\subsubsection*{\textbf{•系统应用边界}}

医学影像辅助诊断系统可以提高病灶筛查和定量测量的效率，但其输出仍会受到影像质量、扫描参数、设备型号、患者配合和训练数据来源等因素影响。呼吸运动可能使胸部CT中的结节边界模糊，金属植入物产生的伪影可能干扰骨折或肿瘤识别，眼底图像中的反光和屈光间质混浊也可能遮挡微小病灶。模型在这些情况下给出的定位框、分割轮廓或概率评分，只能反映当前图像条件下的计算结果，不能直接等同于临床诊断。

不同医院之间的设备、检查协议和患者构成存在差异，模型在开发数据中获得的良好表现未必能够直接延续到新的临床环境。使用薄层CT训练的肺结节模型面对较厚切片时，可能漏掉微小磨玻璃结节；主要使用成年人数据训练的模型，也可能不适用于儿童影像。因此，系统上线前需要在目标医院和目标人群中进行独立测试与外部验证，检查敏感度、特异度、误报率和分割误差是否符合实际应用要求。

系统投入使用后，还需要持续记录模型漏报、误报、医生修正和患者检查条件等信息。当医院更换扫描设备、调整成像协议或更新图像处理软件时，原有数据分布可能发生变化，系统性能应重新评价。医生在使用过程中需要检查模型提示是否符合解剖结构和疾病表现，并结合患者症状、病史、实验室检查及病理结果作出综合判断。医学影像辅助诊断系统的合理作用，是为医生提供更及时、更量化和更易复核的信息，而不是将复杂的医学判断交由算法独立完成。

\subsubsection*{\textbf{•小结}}
本节介绍了医学影像辅助诊断系统的形成基础、系统组成、核心技术、临床流程和应用边界。医学影像数字化使图像能够被计算机存储、测量和分析，CADe和CADx推动影像分析从可疑区域提示走向病变性质判断，深度学习进一步提高了分割、检测、分类和定量测量能力。医学影像辅助诊断系统能够帮助医生提高阅片效率、完成定量分析和发现可疑病灶，但其结果仍需结合原始影像、临床资料和医生经验进行复核。

\subsection{辅助诊断结果与临床判读}

医生打开辅助诊断结果时，看到的可能是胸部CT上的肺结节定位框、颅脑影像中的出血区域轮廓，也可能是眼底病变等级、心室容积和射血分数等定量信息。这些结果是模型根据影像特征生成的分析提示，用于帮助医生定位异常、测量病灶和整理判断线索，但不能直接替代临床诊断。肺结节风险评分较高时，医生仍需观察结节密度、毛刺和随访变化；脑出血区域被自动勾画后，还需检查模型是否误将钙化或伪影划入病灶；心功能参数出现异常时，也要结合图像切面、心律和患者症状判断测量结果是否可靠。

辅助诊断结果通常包括病灶定位、区域分割、类别判断、定量测量和风险提示，不同形式对应不同的临床判读要求。定位结果需要确认模型标记的位置是否符合解剖结构，分割结果需要检查边界是否存在遗漏或过度勾画，分类和风险结果则需要结合模型置信程度及患者实际情况进行解释。眼底筛查系统提示重度病变时，医生需要查看出血、渗出和新生血管等具体依据；肿瘤随访中，病灶体积发生变化还应排除扫描层厚和分割偏差造成的影响。如图\ref{fig:ch9_14}所示，模型输出只有与原始影像、既往检查和临床资料相互对应，并经过医生复核后，才能转化为具有实际意义的诊疗信息。


\par
\noindent\begin{minipage}{\textwidth}
\centering
\BookFigureImage{./images_chapter_07-08/9.14.png}
\captionof{figure}{辅助诊断结果类型与临床判读关系示意图}
\label{fig:ch9_14}
\end{minipage}
\par
\subsubsection*{\textbf{•辅助诊断结果的形成}}

医学影像辅助诊断结果通常经过图像质量检查、目标区域识别、病灶特征分析和结果生成等步骤。系统首先检查影像是否完整、清晰，并判断其是否满足分析要求，随后根据具体任务对肺野、心腔、血管或其他目标区域进行处理。胸部CT中的肺结节分析需要比较可疑阴影在连续切片中的位置、形态和大小，脑出血分析侧重识别异常密度区域，眼底影像分析则关注微血管瘤、出血和渗出等细小病变。不同任务所关注的组织结构和异常特征不同，采用的分析流程也会有所差异。

模型完成分析后，需要将计算结果转换为医生能够查看和核对的信息。检测模型用于标明病灶位置，分割模型描绘目标边界，分类模型给出病变类别及相应评分，定量模型则计算病灶直径、体积或心功能参数。心脏影像系统计算射血分数时，需要先识别舒张末期和收缩末期的心腔轮廓；如果轮廓分割不准确，后续参数也会产生偏差。因此，辅助诊断结果应与原始影像对应显示，使医生能够检查病灶位置、边界和测量依据。

\subsubsection*{\textbf{•结果类型与表达方式}}

医学影像识别结果通常包括定位、分割、分类、定量和提示性信息。定位结果用于标明病灶所在位置，常采用标记框、中心点或热力图等形式呈现；分割结果则描绘病灶的边界和范围，通常以轮廓线或彩色区域叠加在原始影像上。肺结节分析可以先通过定位结果提示可疑结节位置，再利用分割结果计算其直径、体积和形态变化；脑出血分析需要较完整地勾画血肿区域，以便比较复查影像中出血范围和体积的变化。

分类结果用于判断病变所属类别或严重程度，定量结果提供体积、面积、厚度或功能参数等客观数值，提示性结果则根据分析结果给出复查、转诊或进一步检查建议。皮肤镜分析系统可以输出病变类别及相应评分，眼底影像系统可以给出病变分级和转诊提示，心脏影像系统则可提供心室容积、室壁厚度和射血分数等参数。模型评分较高仅说明当前影像与某类病变特征较为接近，不能直接替代临床诊断，因此结果界面还应同时显示病灶位置、分割范围和测量依据，帮助医生核对模型判断。


\subsubsection*{\textbf{•临床判读的基本要求}}

医生判读模型结果时，首先需要核对提示区域是否符合正常解剖结构和相应疾病的影像表现。肺结节标记应结合相邻切片连续观察，避免将圆形血管截面误判为结节；眼底系统标出的红色区域需要区分真实出血、成像反光及其他伪影；脑出血分割结果则要检查是否误将颅内钙化或高密度骨性结构划入病灶范围。只有定位、分割和测量结果能够在原始影像中得到验证，后续分类和定量结果才具有临床参考价值。

模型输出还需要结合患者的年龄、病史、症状和其他检查结果进行解释。相同大小的肺结节出现在长期吸烟的老年患者和无相关危险因素的年轻患者中，临床风险并不相同；PET显示的异常摄取既可能与肿瘤有关，也可能由炎症、感染或正常生理活动引起；心脏射血分数下降还需结合心律、临床症状和心电图结果判断。医生应综合既往影像、实验室检查和随访变化，对模型提示作出接受、修正或排除的决定。

\subsubsection*{\textbf{•结果可靠性与使用边界}}

医学影像的采集质量会直接影响模型分析结果。胸部CT层厚过大可能遗漏微小磨玻璃结节，患者运动可能导致MRI图像模糊，眼底反光会遮挡血管和病变区域，病理切片折叠则可能改变组织结构和细胞形态。当图像质量较差、病灶表现少见，或模型提示与医生观察不一致时，应返回原始影像核查分析区域、病灶边界和关键结构，不能直接采用自动输出结果。

涉及恶性肿瘤、急性脑出血、放疗靶区或复杂心脏异常等高风险情况时，需要通过人工复核、重复测量、专科会诊或多学科讨论进一步确认。系统还应保存原始影像、模型版本、输出结果及医生修改记录，便于追溯结果形成和修正过程。辅助诊断系统适合用于病灶提示、病例排序和重复测量，疾病定性、治疗方案选择及最终诊断仍需由医生结合完整临床信息作出。



\subsubsection*{\textbf{•小结}}

本节介绍了辅助诊断结果的形成过程、主要表达方式、临床判读要求和使用边界。模型可以通过定位、分割、分类、定量测量和风险提示等形式呈现分析结果，帮助医生发现异常、测量病灶和整理诊断线索；这些结果仍需与原始影像、检查参数、患者病史和其他临床资料结合，并经过医生复核后才能用于后续诊疗判断。



\subsection*{本节小结}

本节介绍了医学影像识别与辅助诊断的基本方法与应用。X线、CT、MRI、超声、核医学和数字病理图像在结构显示、功能观察和病灶分析上各有侧重；医学影像人工智能一般经过图像预处理、区域提取、图像分割、特征表示、模型训练和性能评价等环节完成分析。在肺结节检测、眼底病变筛查、皮肤病变识别、病理肿瘤分级和心脏影像分析等任务中，可根据影像特点选择相应的网络结构与方法。辅助诊断系统的输出包括定位、分割、分类、定量测量和风险提示等形式，但其结果会受到图像质量、设备差异和病灶形态的影响，只能作为诊断参考，最终判断仍需由医生结合患者临床资料完成。

\subsection*{本节习题}
\begin{enumerate}
\item 简述X线、CT、MRI和超声在成像原理与临床适用场景上的主要区别。
\item 医学影像识别的一般流程包含哪些环节？图像分割在其中起什么作用？
\item 以肺结节检测为例，说明检测、分割与分类三类任务的输出有何不同。
\item 辅助诊断系统的结果为什么不能直接等同于临床诊断结论？请从图像质量、设备差异和临床背景三个方面说明。
\end{enumerate}

\clearpage
\section{临床决策分析与健康管理}

\begin{sectionkeypoints}
\item 理解临床决策分析与健康管理的基本含义，认识其与单次诊断判断的区别，明确多源数据和连续随访在临床分析中的作用。
\item 掌握临床风险预测与疗效评估的基本思路，理解风险预测关注未来结局，疗效评估关注干预后的状态变化。
\item 熟悉临床决策支持与诊疗路径分析的常见方法，包括规则推理、知识图谱、机器学习模型和推荐模型等。
\item 理解个体诊疗管理中的患者分层、方案支持、过程监测和反馈评价，认识个体差异对诊疗管理策略的影响。
\item 认识健康干预中的连续监测、异常预警和动态调整要求，理解人工智能结果在长期健康管理中的辅助作用和使用边界。
\end{sectionkeypoints}

疾病诊断只是诊疗过程中的一个阶段，医生还需要根据病情变化判断治疗是否有效、风险是否升高以及后续是否需要调整方案。临床决策分析综合病历、检验、影像、用药和生理信号，为风险分层、疗效判断和方案选择提供依据；健康管理则进一步利用随访记录、家庭监测和生活状态，持续观察患者在院外的健康变化。心力衰竭患者出现体重增加、血压波动和气促加重时，这些信息可以共同提示病情恶化，而不是等待下一次住院后才发现问题。

人工智能能够整合不同时间和来源的数据，识别单项检查中不易发现的变化趋势。糖尿病管理可结合连续血糖、用药记录、肾功能和眼底检查调整随访重点，肿瘤治疗可联合影像变化、肿瘤标志物和症状记录评价治疗反应，慢性心血管疾病则可根据家庭血压、心率和既往事件提示风险。如图\ref{fig:ch9_15}所示，多源数据经过整理和分析后形成风险提示、疗效评价与管理建议，再由医生结合患者实际情况决定是否调整检查、治疗和随访安排。

\par
\noindent\begin{minipage}{\textwidth}
\centering
\BookFigureImage{./images_chapter_07-08/9.15.png}
\captionof{figure}{临床决策分析与健康管理闭环示意图}
\label{fig:ch9_15}
\end{minipage}
\par

\subsection{临床风险预测与疗效评估}

患者完成检查或治疗后，医生既要判断未来是否可能出现病情进展、并发症或再次入院，也要观察当前干预是否产生了预期效果。冠心病患者可根据症状、检验指标和既往病史评估心血管事件风险，肿瘤患者则可通过治疗前后的病灶体积、代谢活性和肿瘤标志物变化判断治疗反应。风险预测面向尚未发生的临床结局，疗效评估关注治疗前后已经出现的状态变化，二者共同为随访安排和方案调整提供依据。

风险预测与疗效评估需要整合病史、检验、影像、用药和随访等多时点数据，并保证所用变量与分析目标相对应。心力衰竭患者的再入院风险可结合心功能、肾功能和症状变化进行分析，糖尿病患者的治疗效果则可通过血糖、糖化血红蛋白和并发症指标进行评价。如图\ref{fig:ch9_16}所示，分析过程通常包括目标设定、变量整理、模型训练、性能评价和结果反馈，模型输出还需结合患者个体差异和临床标准进行解释。



\par
\noindent\begin{minipage}{\textwidth}
\centering
\BookFigureImage{./images_chapter_07-08/9.16.png}
\captionof{figure}{临床风险预测与疗效评估流程示意图}
\label{fig:ch9_16}
\end{minipage}
\par

\subsubsection*{\textbf{•风险预测的含义}}

心力衰竭患者出院后，医生可能需要判断其30天内再次入院的可能性；肿瘤患者完成治疗后，则需要关注未来数年内复发或转移的风险。风险预测将患者当前状态与未来临床结局联系起来，预测对象可以是疾病发生、病情进展、并发症、再入院、复发或死亡，结果通常以风险概率、风险等级或预警提示呈现。

预测心力衰竭患者未来30天是否再次入院时，需要将“再入院”设为预测事件，将出院后的30天设为预测时间窗，并以患者出院时作为信息截止时间。模型只能使用此前获得的病史、检验、影像、用药和生理监测信息，不能将再次入院后产生的资料提前纳入。糖尿病并发症风险可结合血糖控制、病程和肾功能变化，肺癌复发风险则可综合肿瘤分期、病理特征和治疗记录。

模型提示高风险并不表示相应事件必然发生，而是说明患者当前特征与既往高风险人群较为相似。医生可以据此调整随访频率、补充检查或提前采取干预措施，例如对再入院风险较高的心力衰竭患者加强体重、血压和气促症状监测，但具体处理仍需结合患者病情和医疗条件确定。


\subsubsection*{\textbf{•预测变量与数据处理}}

预测变量需要围绕临床结局选择，并兼顾医学关联和记录稳定性。心血管事件预测可纳入年龄、血压、血脂、心电图和既往病史，肿瘤复发预测可结合病灶分期、影像变化、病理特征和治疗方式，重症患者风险分析则更关注血压、心率、血氧和实验室指标的连续波动。与结局关系较弱或缺失严重的变量过多时，反而可能掩盖真正有价值的信息。

样本构建通常需要明确纳入标准、输入期和结局观察期，并统一检验单位、时间记录和缺失值处理方式。预测患者出院后是否再次住院时，首次住院期间的资料可以作为输入，再次入院后产生的数据则不能用于预测。连续监测数据还需按照时间顺序排列，使模型能够区分指标持续升高、短暂波动和突然恶化等不同变化。

不良事件比例较低时，模型容易偏向预测人数更多的无事件组，可以通过类别权重、重采样或预测阈值调整减轻这种偏差。随访结束时，部分患者可能尚未发生结局，也无法确定此后是否发生，这种情况称为删失数据，需要使用生存分析方法保留已有随访信息，而不能简单归入未发生事件组。

\subsubsection*{\textbf{•模型构建与性能评价}}

出院随访名单中，有些心力衰竭患者可能在数周内再次住院，有些患者则能长期保持稳定；肿瘤患者出现复发的时间也可能相差数月甚至数年。模型需要根据预测目标和数据特点选择相应方法：固定时间窗内是否发生再入院等二分类结局可采用逻辑回归，复发、并发症或死亡等包含事件发生时间的任务可使用Cox比例风险模型。病史、检验和用药记录等结构化数据还可使用随机森林、梯度提升决策树和XGBoost（extreme gradient boosting，极端梯度提升），连续血糖、心率及多次随访记录则可通过LSTM或Transformer分析动态变化。

数据划分应以患者为单位，避免同一患者的多次住院记录、影像或检验结果分别进入训练集和测试集。若某位患者的首次住院资料用于训练，后续住院资料又用于测试，模型可能通过识别个体特征获得虚高结果。样本数量有限时，可以采用CV比较模型在不同数据划分下的表现，同时检查缺失数据和患者筛选是否造成系统偏差。

固定时间窗内的风险分类可使用敏感度、特异度、F1和AUC评价模型性能。敏感度表示实际发生结局的患者中被模型正确识别为高风险的比例，数值越高，漏掉高风险患者的情况越少；特异度表示实际未发生结局的患者中被正确识别为低风险的比例，数值越高，无效预警越少。F1综合反映模型识别高风险患者的准确性和完整性，数值越高，说明两者之间的平衡越好；AUC反映模型整体区分高风险与低风险患者的能力，数值越接近1，区分能力越强。对于包含随访时间的生存预测，C-index（concordance index，一致性指数）用于评价风险排序是否正确，数值越高，说明模型越能将较早发生结局的患者排在较高风险位置。

在评价模型预测概率的可靠性时，可以将模型给出的风险水平与患者的实际结局进行比较。若模型将100名患者的风险均预测为20\%，实际发生相应结局的人数也应接近20人；若实际仅有5人，则说明模型高估了风险。将预测结果转化为临床预警时，还需根据疾病严重程度、干预成本和漏诊后果设置阈值，因为阈值过低会使较多低风险患者触发提醒，增加医务人员的核查负担，阈值过高则可能遗漏真正需要提前干预的患者。因此，模型性能评价需要综合考察区分能力、概率可靠性和不同阈值下的临床应用后果。



\subsubsection*{\textbf{•疗效评估方法}}

肿瘤患者接受治疗后，医生需要结合病灶大小、代谢活性和症状变化判断治疗反应；心力衰竭患者则可通过运动耐量、心室功能和再入院情况评价病情改善程度。疗效评估应根据治疗目标选择观察指标，并设置治疗前、治疗中、治疗后和随访等时间点，以分析改善幅度、持续时间及患者之间的反应差异。

单项指标通常难以完整反映治疗效果，因此可以整合影像、检验、生理信号和功能评分进行综合评价。肿瘤体积缩小但PET显示代谢活性仍较高时，需要联合结构与功能信息判断残余病灶；康复患者的影像变化不明显时，还可结合步行能力、肌力和日常活动记录评价功能恢复。模型也可分析不同疗效患者的共同特征，为后续治疗方案调整提供参考。

比较不同治疗方案时，患者在年龄、病情和合并症方面的基础差异可能影响结果。PSM（propensity score matching，倾向评分匹配）通过匹配基础特征相近的患者提高组间可比性，IPW（inverse probability weighting，逆概率加权）则根据患者接受不同治疗的概率调整样本权重，因果森林可用于分析不同患者之间的治疗效果差异。这些方法能够减轻部分混杂因素的影响，但分析结果仍需结合研究设计、数据质量和临床证据进行判断。
\subsubsection*{\textbf{•小结}}

本节介绍了临床风险预测与疗效评估的基本任务、数据处理、模型构建和结果评价方法。风险预测用于估计疾病进展和不良事件发生的可能性，疗效评估用于分析治疗前后的状态变化与个体反应，常用方法包括分类模型、生存分析、时序模型和因果推断。相关结果可为风险分层、治疗调整和后续健康管理提供参考。


\subsection{临床决策支持与诊疗路径}


一名老年患者因胸闷、气短进入急诊时，心力衰竭、急性冠脉综合征和肺部感染都可能引起相似症状，医生需要在较短时间内综合既往病史、心电图、肌钙蛋白、脑钠肽和胸部影像等资料。不同检查提供的线索可能相互补充，也可能存在矛盾，临床决策分析可以按照信息产生的时间和医学关联整理关键变化，帮助医生确定需要优先排查的疾病与危险因素。

CDSS（clinical decision support system，临床决策支持系统）综合患者数据、医学知识和计算模型，为诊断鉴别、检查选择、用药安全和治疗方案提供提示。肌钙蛋白升高并伴有缺血性心电图改变时，系统可提示进一步评估急性冠脉综合征；患者肾功能下降且正在使用经肾脏排泄的药物时，系统可提醒医生核对剂量。系统提供的结果主要用于集中呈现关键信息和降低遗漏风险，最终诊断仍需结合症状变化、完整检查结果及医生经验作出。

诊疗路径按照疾病进程组织检查、治疗、观察和随访环节，CDSS可以根据患者当前状态提示下一步需要完成的任务。如图\ref{fig:ch9_17}所示，患者资料与医学知识库、诊疗指南共同形成诊断提示、风险提醒和处理建议，再由医生判断是否采纳。病情复杂、证据不一致或患者存在多种合并症时，既定路径需要根据个体情况调整，不能机械替代临床判断。


\par
\noindent\begin{minipage}{\textwidth}
\centering
\BookFigureImage{./images_chapter_07-08/9.17.png}
\captionof{figure}{临床决策支持系统与诊疗路径示意图}
\label{fig:ch9_17}
\end{minipage}
\par

\subsubsection*{\textbf{•临床决策支持}}

医生在诊疗过程中需要同时核对症状、检查结果、既往疾病、当前用药和诊疗指南，信息数量较多时容易出现遗漏。CDSS可以将患者资料与医学知识进行关联，为诊断鉴别、检查选择和用药安全提供提示。检验结果显示血钾明显升高时，系统可提醒医生及时复查并评估心律失常风险；患者同时使用抗凝药和可能增加出血风险的药物时，系统也可提示核查药物相互作用。

危急值、重复用药、药物相互作用和诊疗流程等要求较为明确，通常可以转化为规则库中的判断条件，使系统在满足特定条件时自动发出提醒。规则提示具有清晰、易解释的特点，但难以覆盖多病共存、少见病例和证据相互矛盾的情况，因此系统给出的建议仍需结合患者完整资料、诊疗指南和医生判断进行分析。


\subsubsection*{\textbf{•知识图谱}}

疾病、症状、检查、药物和治疗方案之间存在复杂联系，医学知识图谱通过节点表示这些医学对象，并利用关系描述疾病表现、检查依据、药物作用和治疗限制。肿瘤诊疗中，病理类型、基因变异、药物靶点和治疗方案可以形成关联网络，系统能够根据患者的病理报告和基因检测结果查找相关治疗证据，帮助医生整理分散在不同资料中的诊疗线索。

感染性疾病的分析还可将感染部位、病原体、药敏结果和抗菌药物联系起来，当影像和炎症指标提示感染时，系统可以提示需要完善的病原学检查，并核对药物选择与药敏结果是否一致。知识图谱中的指南、药品和疾病知识需要随临床证据及时更新，若关系缺失或信息过期，系统生成的建议可能不再适用于当前诊疗要求，因此仍需由医生结合患者资料进行核查。

\subsubsection*{\textbf{•诊疗路径}}

诊疗路径按照疾病进程组织评估、检查、诊断、治疗和随访环节，使医务人员能够根据患者当前状态确定下一步任务。急性缺血性脑卒中的救治对时间要求较高，系统可结合发病时间、颅脑影像、凝血功能和用药情况，提示完成溶栓或取栓评估；患者进入恢复期后，路径还可继续安排功能评价、康复训练和复发风险管理，从而保持院内治疗与出院随访的连续性。

标准化流程有助于减少检查遗漏、重复操作和处置延迟，但患者出现严重合并症、治疗禁忌或不典型表现时，原有步骤可能不再适用。系统可以提示路径偏离及其原因，医生则需结合病情变化、患者意愿和医疗条件调整检查与治疗安排，使诊疗路径成为临床工作的参考，而不是限制个体化决策的固定程序。
\subsubsection*{\textbf{•决策推理}}

临床决策支持需要先从病历、检查报告和出院记录中提取症状、诊断、药物及时间信息，再通过规则推理核对危急值、用药禁忌和诊疗要求。机器学习可利用历史病例评估疾病风险与治疗反应，知识图谱则将患者资料与疾病、检查、药物和治疗知识关联起来，使分散的信息能够形成可供判断的诊疗线索。
手术前评估和术后管理中，系统可以利用NLP从病历中提取既往疾病、过敏史和当前用药，通过规则核对禁食要求、检查项目及抗凝药物是否需要调整，并结合预测模型评估术后感染、出血或其他并发症风险。系统输出还应呈现触发提示的患者数据和医学依据，使医生能够核查信息来源，并判断相关建议是否适用于当前患者。



\subsubsection*{\textbf{•应用边界}}

当关键检查尚未完成、病历记录存在缺失或知识库更新不及时，系统可能无法获得充分证据，生成的建议也可能出现偏差。罕见病、多病共存和异常治疗反应往往超出既有规则与训练数据的覆盖范围，因此急危重症和高风险治疗仍需结合现场观察、多学科意见、患者意愿及最新临床证据作出决定。

系统频繁发出缺乏针对性的提示时，医务人员容易产生提醒疲劳，进而忽视真正需要立即处理的信息。提示设计应优先呈现严重药物相互作用、危急检验结果和病情急剧恶化等高风险事件，并提供明确的处置依据；低风险信息可采用集中展示或分级提醒，减少无效干扰，避免增加临床信息负担。


\subsubsection*{\textbf{•小结}}
本节介绍了临床决策支持与诊疗路径的基本内容，临床决策支持系统通过规则、知识图谱和智能模型整合患者资料，为诊断、检查、用药和治疗提供提示；诊疗路径则把疾病诊断、治疗和随访管理组织为相对清晰的流程，但系统结果仍受数据质量和知识更新等因素影响，需要由医生结合患者实际情况进行判断。

\subsection{个体诊疗管理与健康干预}

家庭监测中的血压、血糖或心率可能随用药、饮食、运动和睡眠状态不断变化，单次门诊检查难以完整反映患者的长期健康状况。个体诊疗管理需要综合病史、检验、影像、生理监测、用药和随访资料，分析患者的风险水平、治疗反应和依从性，并据此调整复查频率、治疗方案与生活方式干预。

人工智能可以整理持续积累的健康数据，识别指标异常和病情变化趋势，辅助患者分层与干预方案匹配。如图\ref{fig:ch9_18}所示，患者资料经过状态评估后形成诊疗和健康干预建议，后续监测结果再反馈至分析环节，使管理方案能够根据治疗效果和健康状态持续调整。


\par
\noindent\begin{minipage}{\textwidth}
\centering
\BookFigureImage{./images_chapter_07-08/9.18.png}
\captionof{figure}{个体诊疗管理与健康干预闭环示意图}
\label{fig:ch9_18}
\end{minipage}
\par

\subsubsection*{\textbf{•个体诊疗管理的含义}}

治疗方案实施后，部分患者的病情能够较快改善，另一些患者可能因合并症、药物耐受性或生活习惯不同而需要调整方案。个体诊疗管理根据疾病状态、身体条件、既往治疗反应和生活方式等信息，持续分析患者在诊断、治疗和随访过程中的变化，使治疗选择与后续管理更加符合个体需求。

人工智能开展个体诊疗分析时，通常需要整合病历、检验、影像、用药、监测和随访数据，形成包含疾病特征、风险水平、治疗反应、依从性和随访状态的患者画像。系统可在此基础上进行患者分层、相似病例匹配和方案辅助分析，帮助医生识别需要重点监测或及时调整治疗的患者。

临床研究提供的证据主要反映群体规律，而具体患者的治疗效果可能受到多种因素共同影响。人工智能可以分析历史病例中患者特征、风险水平与治疗反应之间的关系，将群体数据中的规律转化为个体诊疗参考，但相关结果仍需结合患者当前状态和临床证据进行判断。


\subsubsection*{\textbf{•个体差异与数据基础}}

患者的年龄、性别、体重、基础疾病和遗传背景相对稳定，而生命体征、检验指标、影像表现、症状变化和用药反应会随时间改变，这些信息共同构成个体差异的数据基础。结构化临床数据用于描述基本状态和指标水平，医学影像反映组织结构与病灶变化，生理信号记录连续状态，随访数据则呈现治疗反应和长期管理过程。

在分析个体差异时，需要将分散在病历、检查、监测和随访系统中的数据进行整合，并通过变量筛选和特征表示提取与疾病状态、风险水平及治疗反应相关的信息。统计分析、聚类分析、降维、特征选择、表示学习和多模态融合等方法，可将不同来源和形式的数据转化为模型能够处理的特征，同时减少冗余信息对分析结果的影响。

患者分层可以利用聚类分析、决策树、随机森林、梯度提升决策树或深度学习模型，将患者划分为不同风险层级或反应类型。分层结果可用于识别需要密切监测的患者、可能对特定干预更敏感的人群以及依从性较低的群体，为随访频率、干预重点和管理方式的差异化调整提供参考。


\subsubsection*{\textbf{•个体化方案制定}}

在个体化方案制定过程中，患者的疾病状态、风险水平、治疗反应、身体条件和生活习惯需要转化为具体的诊疗与健康管理目标。方案内容可以包括药物选择、治疗强度、检查项目、随访频率以及饮食、运动和康复安排。年龄较大且存在多种合并症的高血压患者，治疗目标和用药强度需要兼顾降压效果与不良反应风险，不能直接套用年轻患者的管理方式。

人工智能可以综合患者当前资料、既往治疗记录和相似病例的管理结果，比较不同候选方案可能带来的获益与风险，并提示需要重点关注的药物反应、并发症和随访指标。输出结果应说明推荐依据、适用条件和潜在风险，使医生能够结合临床指南、患者意愿和实际医疗条件进行选择，而不是仅给出缺少解释的方案排序。

在持续监测过程中，检验指标、症状、生理信号和依从性变化可以反映方案是否有效。治疗效果稳定且未出现明显不良反应时，可维持原有管理安排；病情恶化、指标异常或患者难以长期执行时，则需要调整治疗强度、随访频率或生活方式干预。通过评估、干预、监测和反馈的连续循环，个体诊疗方案能够随患者状态变化进行更新。

\subsubsection*{\textbf{•健康干预与过程监测}}

患者离院后，用药、饮食、运动、康复训练和定期复查仍会持续影响疾病控制效果。健康干预将诊疗管理延伸到日常生活，通过观察健康指标和行为执行情况，及时发现病情变化或方案落实中的问题。高血压患者可持续记录家庭血压和用药情况，术后康复患者则可结合活动量、训练完成度和疼痛变化评估恢复进程。

可穿戴设备、移动健康平台和远程监测系统能够连续采集心率、血压、血糖、活动量及睡眠等信息，人工智能可据此分析变化趋势、识别异常波动并形成风险提示。当患者连续多日血糖偏高或康复训练完成度下降时，系统可以提醒其复测指标、调整行为安排或及时联系医务人员，并根据患者状态提供复查、用药和生活方式方面的个体化提示。

连续监测数据容易受到设备误差、记录缺失、测量条件和个体基线差异的影响，因此异常结果需要结合变化持续时间、既往水平和患者症状进行判断。稳定且可追踪的监测记录能够反映干预效果，并为后续调整治疗强度、随访频率和健康管理方案提供依据。


\subsubsection*{\textbf{•干预效果评价与管理更新}}

在持续干预过程中，患者的症状、检验指标、功能状态和生活质量需要与干预前水平进行比较，以判断治疗目标是否达到、改善效果能否维持以及是否出现新的风险。慢性阻塞性肺疾病患者接受康复训练后，可结合呼吸困难程度、运动耐量和急性加重次数评价干预效果，避免仅凭单次检查作出判断。

反馈评价还需关注患者对方案的执行情况。用药记录、随访结果和设备监测数据能够反映服药、运动及康复训练是否按计划完成，并帮助区分干预效果不佳是由方案本身、依从性不足还是病情变化所致。人工智能可以整合最新数据，更新患者画像、风险等级和治疗反应，为后续管理提供依据。

当风险水平升高、症状加重或干预目标未达到时，医生可据此调整治疗强度、监测指标和随访频率；患者状态稳定且方案执行良好时，则可维持当前安排。通过持续评价和信息反馈，个体诊疗管理能够随患者状态变化不断更新。


\subsubsection*{\textbf{•小结}}

本节介绍了个体诊疗管理与健康干预的基本过程，包括患者状态分析、个体化方案制定、院外监测和干预效果评价。人工智能可以整合多源数据，辅助患者分层、异常识别和管理方案更新，但相关结果仍需结合患者实际情况和临床判断使用。




\subsection*{本节小结}

本节介绍了临床决策分析与健康管理的基本思路。与单次诊断不同，临床决策分析综合病历、检验、影像、用药和生理信号等多源数据，用于风险分层、疗效评估和方案选择；健康管理则借助随访记录和院外监测，持续观察患者的长期健康变化。风险预测关注未来结局，疗效评估关注干预后的状态变化，临床决策支持可通过规则推理、知识图谱和机器学习模型等方式为诊疗路径提供参考。在个体诊疗管理中，患者分层、方案支持、过程监测和反馈评价共同构成动态闭环。人工智能有助于发现单项检查中不易察觉的变化趋势，但其结果仍需结合患者实际情况，由医生判断是否调整检查、治疗和随访安排。

\subsection*{本节习题}
\begin{enumerate}
\item 临床风险预测与疗效评估在关注点上有何区别？请各举一个应用例子。
\item 临床决策支持常用哪些方法？知识图谱与机器学习模型在其中各有什么特点？
\item 个体诊疗管理为什么强调患者分层与过程监测？试结合慢性病管理加以说明。
\item 在长期健康管理中，连续监测数据可能出现设备误差或数据缺失，应如何看待人工智能给出的预警结果？
\end{enumerate}

% ============================================================
% 全章回顾：医学人工智能的四大支柱
% ============================================================
\subsection*{全章回顾：医学人工智能的四大支柱}

电子病历、检验指标、医学影像、生理信号、组学数据——一名患者从挂号就诊到出院随访，这五类数据不断产生、不断积累。它们来源不同、格式不同、医学含义也不同，但都指向同一个目标：理解疾病、评估风险、辅助决策。这一章就沿着"数据→影像→决策→管理"这条主线，四个支柱一个一个走了一遍：

\begin{enumerate}
\def\labelenumi{\arabic{enumi}.}
\setlength{\itemsep}{3pt}
\item \textbf{医学数据处理与任务类型。}电子病历、检验指标、医学影像、生理信号和组学数据是医学人工智能的五类基础数据。它们的共同特征是多源异构、专业性强、动态变化和敏感性高。数据采集、清洗、标注、标准化和多源融合——这五个处理环节决定了模型能学到什么、学不到什么。分类、检测、分割、预测和生成是五类典型任务，不同任务对数据格式和标注方式的要求截然不同。
\item \textbf{医学影像识别与辅助诊断。}从X射线到CT、MRI再到数字病理全切片，影像类型不同，分析难度也不同。人工智能通过图像预处理、病灶定位与分割、特征提取和模型识别，从影像中提取结构与病变信息，并以定位框、分割轮廓、疾病类别和风险提示等形式为医生提供参考。模型输出的是辅助信息而非最终诊断——医生仍需结合病史、体征和其他检查结果进行综合判断。
\item \textbf{临床决策分析与健康管理。}临床决策不止于单次诊断。风险预测关注"未来会发生什么"，疗效评估关注"干预后改变了什么"，临床决策支持则通过规则推理、知识图谱和机器学习模型为诊疗路径提供参考。三者共同将医学人工智能从"看图说话"推向"推理辅助"。
\item \textbf{个体诊疗管理与健康干预。}同样的疾病、不同的患者——年龄、合并症、治疗反应和生活习惯的差异决定了管理方案必须个体化。患者分层、方案支持、过程监测和反馈评价构成动态闭环，人工智能帮助识别异常趋势和风险变化，但最终调整检查、治疗和随访安排仍需医生结合患者实际情况做出判断。
\end{enumerate}

四个支柱走完，就一句话：\textbf{医学人工智能的价值在于提高信息处理效率和发现数据中的异常线索，但模型输出不能脱离原始数据和患者临床背景进行解释。}它能帮助医生看到更多、算得更快，但"看到什么、意味着什么"——这个判断只能由医生做出。


% ============================================================
% 章末习题
% ============================================================
\subsection*{章末习题}

\subsubsection*{一、简答题}

\begin{enumerate}
\def\labelenumi{\arabic{enumi}.}
\setlength{\itemsep}{6pt}

\item \textbf{医学数据的类型与特征。}医学人工智能常用的数据类型有哪些？列举至少五种，并简述其主要记录内容和医学用途。这些数据在结构形式、时间尺度和信息密度方面有哪些基本特征？请结合具体诊疗场景加以说明。

\item \textbf{医学人工智能的典型任务。}分类、检测、分割和风险预测是医学人工智能的四种典型任务。用自己的话解释每种任务的基本作用，并分别举出一个实际的医学应用场景。

\item \textbf{医学影像类型与处理流程。}常见医学影像类型有哪些？简述X射线成像、CT、MRI、超声成像和核医学成像的主要特点。在此基础上，说明医学影像人工智能分析通常包括哪些基本环节——按处理顺序解释图像预处理、病灶定位、图像分割、特征分析和结果输出各自的作用。

\item \textbf{AI辅助诊断结果的解读边界。}人工智能辅助诊断系统输出了病灶位置、分割范围、疾病类别和风险提示。结合这些输出结果，说明它们能为医生提供哪些信息，以及医生在使用这些结果时还需要核对哪些临床资料。为什么模型输出不能直接等同于临床诊断结论？

\item \textbf{临床决策支持中的人机分工。}临床决策支持系统发现患者存在药物相互作用风险，并提示医生核对治疗方案。说明规则库、知识图谱、患者资料和诊疗路径在这一过程中分别发挥什么作用，并分析系统提示不能直接替代医生决策的原因。

\end{enumerate}

\subsubsection*{二、操作题}

\begin{enumerate}
\def\labelenumi{\arabic{enumi}.}
\setlength{\itemsep}{6pt}
\setcounter{enumi}{5}

\item \textbf{数字病理图像分析方案设计。}数字病理图像具有图像尺寸大、组织结构复杂和细胞数量多等特点。请选择一个你了解的病理场景（如肿瘤分级、细胞计数或组织分类），设计一个完整的人工智能辅助分析方案，包括：
  \begin{enumerate}
    \item 图像预处理方法的选择及其理由
    \item 标注策略的制定——标注什么、需要多少样本、谁来标注
    \item 适用的模型类型及其原因
    \item 评估指标的选取——为什么选择这些指标？
  \end{enumerate}

\item \textbf{再住院风险预测模型设计。}某医院计划预测心力衰竭患者出院后30天内再次住院的风险。请完成以下任务：
  \begin{enumerate}
    \item 确定预测目标、输入数据范围（至少列出5类变量）和观察时间窗口
    \item 说明训练集与测试集的划分方法，并分析划分不当可能造成什么问题
    \item 选择合适的评估指标并说明理由——为什么准确率（accuracy）可能不是最好的选择？
  \end{enumerate}

\item \textbf{疗效评估的数据选择与分析。}肿瘤患者接受治疗后，医生需要判断病灶是否缩小、代谢活性是否下降以及症状是否改善。请完成：
  \begin{enumerate}
    \item 列出至少三类用于疗效评估的数据来源
    \item 说明为什么不能只依据单项指标（如肿瘤直径）判断治疗效果
    \item 设计一个多指标综合评估框架，说明各类指标的权重分配逻辑
  \end{enumerate}

\end{enumerate}

\subsubsection*{三、综合题}

\begin{enumerate}
\def\labelenumi{\arabic{enumi}.}
\setlength{\itemsep}{6pt}
\setcounter{enumi}{8}

\item \textbf{个体诊疗管理方案设计。}选择高血压、糖尿病、心力衰竭或术后康复中的一种情况，设计一个完整的个体诊疗管理方案。方案中应包括：
  \begin{enumerate}
    \item 患者状态评估的指标体系——需要采集哪些数据、以什么频率采集
    \item 患者分层策略——根据哪些指标将患者分为不同的管理级别
    \item 过程监测方案——院外监测的技术手段、关键预警指标和异常判定标准
    \item 反馈调整机制——在什么条件下调整治疗方案、调整的方向是什么
  \end{enumerate}

\item \textbf{批判性反思：医学AI的边界。}本章从一开始就强调"模型输出不能脱离原始数据和患者临床背景进行解释"。结合你在以上简答题和操作题中的分析，写一段不少于300字的反思，讨论：
  \begin{enumerate}
    \item 在你设计的方案中，AI可以在哪些环节真正提高效率？哪些环节AI的帮助有限甚至可能产生误导？
    \item 数据质量问题（如记录缺失、标注差异、设备变化）可能如何在你的方案中产生连锁反应？
    \item 你如何理解医学人工智能中的"人机协作"——医生的判断在什么情况下不可替代？
  \end{enumerate}

\end{enumerate}



\section*{术语注释表}

\begin{center}
\small
\renewcommand{\arraystretch}{1.18}
\begin{longtable}{p{0.18\textwidth}p{0.45\textwidth}p{0.27\textwidth}}
\toprule
缩写 & 英文全称 & 中文名称\\
\midrule
\endfirsthead
\toprule
缩写 & 英文全称 & 中文名称\\
\midrule
\endhead
ML & machine learning & 机器学习\\
DL & deep learning & 深度学习\\
CT & computed tomography & 计算机体层成像\\
MRI & magnetic resonance imaging & 磁共振成像\\
NLP & natural language processing & 自然语言处理\\
DSA & digital subtraction angiography & 数字减影血管造影\\
PET & positron emission tomography & 正电子发射断层扫描\\
SPECT & single-photon emission computed tomography & 单光子发射计算机体层成像\\
MS & mass spectrometry & 质谱分析\\
DNA & deoxyribonucleic acid & 脱氧核糖核酸\\
WBC & white blood cell count & 白细胞计数\\
ICD & International Classification of Diseases & 国际疾病分类\\
DICOM & Digital Imaging and Communications in Medicine & 医学数字成像与通信\\
CRP & C-reactive protein & C反应蛋白\\
PCA & principal component analysis & 主成分分析\\
CNN & convolutional neural network & 卷积神经网络\\
RNN & recurrent neural network & 循环神经网络\\
LR & logistic regression & 逻辑回归\\
SVM & support vector machine & 支持向量机\\
RF & random forest & 随机森林\\
GBDT & gradient boosting decision tree & 梯度提升决策树\\
LSTM & long short-term memory & 长短期记忆网络\\
KG & knowledge graph & 知识图谱\\
RL & reinforcement learning & 强化学习\\
T1WI & T1-weighted imaging & T1加权成像\\
T2WI & T2-weighted imaging & T2加权成像\\
DWI & diffusion-weighted imaging & 弥散加权成像\\
FLAIR & fluid-attenuated inversion recovery & 液体衰减反转恢复序列\\
FCN & fully convolutional network & 全卷积网络\\
Mask R-CNN & mask region-based convolutional neural network & 掩膜区域卷积神经网络\\
3D CNN & three-dimensional convolutional neural network & 三维卷积神经网络\\
KNN & k-nearest neighbors & K近邻\\
VGG & Visual Geometry Group network & 视觉几何组网络\\
ResNet & residual network & 残差网络\\
DenseNet & densely connected convolutional network & 密集连接卷积网络\\
CV & cross-validation & 交叉验证\\
F1 & F1 score & F1分数\\
AUC & area under the receiver operating characteristic curve & 受试者工作特征曲线下面积\\
DSC & Dice similarity coefficient & Dice相似系数\\
IoU & intersection over union & 交并比\\
LDCT & low-dose computed tomography & 低剂量计算机体层成像\\
DR & diabetic retinopathy & 糖尿病视网膜病变\\
SK & seborrheic keratosis & 脂溢性角化\\
BCC & basal cell carcinoma & 基底细胞癌\\
WSI & whole slide image & 全切片图像\\
MIL & multiple instance learning & 多实例学习\\
EF & ejection fraction & 射血分数\\
CAD & computer-aided diagnosis & 计算机辅助诊断\\
CR & computed radiography & 计算机X射线摄影\\
DR & digital radiography & 数字X射线摄影\\
CADe & computer-aided detection & 计算机辅助检测\\
CADx & computer-aided diagnosis & 计算机辅助诊断\\
PACS & picture archiving and communication system & 影像归档与通信系统\\
XGBoost & extreme gradient boosting & 极端梯度提升\\
C-index & concordance index & 一致性指数\\
PSM & propensity score matching & 倾向评分匹配\\
IPW & inverse probability weighting & 逆概率加权\\
CDSS & clinical decision support system & 临床决策支持系统\\
\bottomrule
\end{longtable}
\end{center}

